\chapter{RUANG LINEAR BERNORMA}

\section{Norma}
Dalam dunia analisis, penghuni yang paling dominan adalah ruang fungsi, yakni ruang vektor
yang terdiri atas fungsi-fungsi bernilai real atau kompleks. Agar menarik bagi seorang analis,
ruang semacam itu harus dilengkapi dengan suatu topologi. Sering kali topologi tersebut
merupakan topologi metrik, yang pada gilirannya kerap berasal dari suatu norma (lihat
proposisi~\ref{0001503}).

\begin{exam}\label{C023125} Untuk $x = (x_1,\dots, x_n) \in \K^n$, misalkan $\norm{x} =
\left(\sum_{k=1}^n {\abs{x_k}}^2\right)^{1/2}$. Ini adalah
 \index{biasa!norma!pada $\K^n$}%
 \index{norma!biasa!pada $\K^n$}%
\df{norma biasa} (atau
 \index{Euklides!norma!pada $\K^n$}%
 \index{norma!Euklides!pada $\K^n$}%
 \index{K@$\K^n$!sebagai ruang linear bernorma}%
 \index{bernorma!ruang linear!$\K^n$ sebagai}%
\df{norma Euklides}) pada $\K^n$; kecuali jika secara eksplisit dinyatakan sebaliknya,
ketika dipandang sebagai ruang linear bernorma, $\K^n$ selalu dianggap memiliki norma ini.
Jelas bahwa norma ini menginduksi metrik biasa (Euklides) pada~$\K^n$.
\end{exam}

\begin{exam}\label{C023127} Untuk $x = (x_1,\dots,x_n) \in \K^n$, misalkan
 \index{<unopn@$\norm{\hphantom{x}}_1$ (norma-$1$)}%
$\norm{x}_1 = \sum_{k=1}^n \abs{x_k}$. Mudah dilihat bahwa fungsi $x \mapsto \norm{x}_1$
merupakan norma pada~$\K^n$. Norma ini kadang-kadang disebut
 \index{norma!$1$- (pada $\K^n$)}%
 \index{satu@norma-$1$!pada $\K^n$}%
 \index{K@$\K^n$!sebagai ruang linear bernorma}%
 \index{bernorma!ruang linear!$\K^n$ sebagai}%
\df{norma-$1$} pada~$\K^n$. Norma ini menginduksi apa yang disebut \emph{metrik taksi}
pada~$\R^2$.
\end{exam}

Dari enam contoh berikutnya, contoh terakhir adalah yang paling umum. Perhatikan bahwa
semua contoh lainnya merupakan kasus khususnya atau subruang dari kasus-kasus khususnya.

\begin{exam}\label{C023129} Untuk $x = (x_1,\dots,x_n) \in \K^n$, misalkan
 \index{<unopn@$\norm{\hphantom{x}}_u$ (norma seragam)}%
 \index{<unopn@$\norm{\hphantom{x}}_\infty$ (norma seragam)}%
$\norm{x}_\infty = \max\{\,\abs{x_k} \colon 1 \le k \le n\}$. Ini mendefinisikan suatu
norma pada $\K^n$; norma tersebut adalah
 \index{seragam!norma!pada~$\K^n$}%
 \index{norma!seragam!pada $\K^n$}%
 \index{K@$\K^n$!sebagai ruang linear bernorma}%
 \index{bernorma!ruang linear!$\K^n$ sebagai}%
\df{norma seragam} pada~$\K^n$. Norma ini menginduksi metrik seragam pada~$\K^n$.
Notasi alternatif untuk $\norm{x}_\infty$ adalah $\norm{x}_U$.
\end{exam}

\begin{exam}\label{C0231301} Himpunan $l_\infty$ dari semua barisan bilangan kompleks yang
terbatas jelas merupakan ruang vektor di bawah operasi penjumlahan dan perkalian skalar
titik demi titik. Untuk $x = (x_1,x_2,\dots) \in l_\infty$, misalkan
 \index{<unopn@$\norm{\hphantom{x}}_\infty$ (norma seragam)}%
$\norm{x}_\infty = \sup\{\,\abs{x_k} \colon 1 \le k\}$. Ini mendefinisikan suatu norma pada
$l_\infty$; norma tersebut adalah
 \index{seragam!norma!pada~$l_\infty$}%
 \index{norma!seragam!pada $l_\infty$}%
 \index{l@$l_\infty$!sebagai ruang linear bernorma}%
 \index{l@$l_\infty$!barisan terbatas}%
 \index{bernorma!ruang linear!$l_\infty$ sebagai}%
\df{norma seragam} pada~$l_\infty$.
\end{exam}

\begin{exam}\label{C0231302} Salah satu subruang dari~\ref{C0231301} adalah ruang $c$ yang
terdiri atas semua
 \index{c@$c$!sebagai ruang linear bernorma}%
 \index{c@$c$!barisan konvergen}%
 \index{bernorma!ruang linear!$c$ sebagai}%
barisan bilangan kompleks yang konvergen. Di bawah norma seragam, ruang ini merupakan
ruang linear bernorma.
\end{exam}

\begin{exam}\label{C0231303} Di dalam $c$ terdapat ruang linear bernorma lain yang menarik,
yaitu $\sbsb c0$
 \index{c@$c_0$!barisan yang konvergen ke nol}%
 \index{c@$c_0$!sebagai ruang linear bernorma}%
 \index{bernorma!ruang linear!$\sbsb c0$ sebagai}%
yakni ruang semua barisan bilangan kompleks yang konvergen ke nol (sekali lagi dengan
norma seragam).
\end{exam}

\begin{exam}\label{C023131} Misalkan $S$ suatu himpunan tak kosong. Jika $f$ merupakan fungsi
bernilai $\K$ yang terbatas pada $S$, misalkan
  \[ \norm f_u := \sup\{\abs{f(x)} \colon x \in S\}\,. \]
Ini adalah norma pada ruang vektor $\fml B(S)$ dari semua fungsi skalar yang
terbatas pada $S$ dan disebut
 \index{<unopn@$\norm{\hphantom{x}}_u$ (norma seragam)}%
 \index{seragam!norma!pada~$\fml B(S)$}%
 \index{norma!seragam!pada~$\fml B(S)$}%
 \index{B@$\fml B(S)$!fungsi terbatas pada~$S$}%
 \index{B@$\fml B(S)$!sebagai ruang linear bernorma}%
 \index{bernorma!ruang linear!$\fml B(S)$ sebagai}%
\df{norma seragam}. Jelas bahwa norma ini menghasilkan metrik seragam pada~$\fml B(S)$.
Perhatikan bahwa contoh~\ref{C023129} dan~\ref{C0231301} merupakan kasus khusus dari
contoh ini. (Ambil $S = \N_n := \{1,2,\dots,n\}$ atau $S = \N$.)
\end{exam}











\begin{exam}\label{C023133} Mudah untuk membuat perumuman yang cukup luas atas contoh
sebelumnya dengan mengganti medan $\K$ dengan ruang linear bernorma sebarang. Maka,
misalkan $S$ suatu himpunan tak kosong dan $V$ suatu ruang linear bernorma. Untuk $f$
dalam ruang vektor $\fml B(S,V)$ dari semua fungsi bernilai $V$ yang terbatas pada~$S$,
misalkan
  \[ \norm f_u := \sup\{\norm{f(x)} \colon x \in S\}\,. \]
Maka ini adalah suatu norma dan disebut
 \index{<unopn@$\norm{\hphantom{x}}_u$ (norma seragam)}%
 \index{seragam!norma!pada~$\fml B(S,V)$}%
 \index{norma!seragam!pada~$\fml B(S,V)$}%
 \index{B@$\fml B(S,V)$!fungsi bernilai $V$ yang terbatas pada~$S$}%
 \index{B@$\fml B(S,V)$!sebagai ruang linear bernorma}%
 \index{bernorma!ruang linear!$\fml B(S,V)$ sebagai}%
\df{norma seragam} pada $\fml B(S,V)$. (Mengapa kata ``mudah'' dalam kalimat pertama
contoh ini tepat digunakan?)
\end{exam}

\begin{defn} Misalkan $S$ suatu himpunan, $V$ suatu ruang linear bernorma, dan
 \index{F@$\fml F(S,V)$ (fungsi bernilai $V$ pada~$S$)}%
$\fml F(S,V)$ ruang vektor dari semua fungsi bernilai $V$ pada~$S$. Perhatikan suatu
barisan $(f_n)$ dari fungsi-fungsi dalam $\fml F(S,V)$. Jika terdapat fungsi $g$ dalam
$\fml F(S,V)$ sedemikian sehingga
   \[ \sup\{\norm{f_n(x) - g(x)}\colon x \in S\} \sto 0
      \quad \text{ketika $n \sto \infty$,} \]
maka kita mengatakan bahwa barisan $(f_n)$
 \index{<arrowto@$f_n \sto g$ (unif) (konvergensi seragam)}%
 \index{seragam!konvergensi}%
 \index{konvergensi!seragam}%
\df{konvergen seragam} ke~$g$ dan menulis $f_n \sto g \text{\,(unif)}$. Fungsi $g$
adalah
 \index{seragam!limit}%
 \index{limit!seragam}%
\df{limit seragam} dari barisan $(f_n)$. Perhatikan bahwa jika $g$ dan semua $f_n$
termasuk dalam $\fml B(S,V)$, maka konvergensi seragam $(f_n)$ ke $g$ tidak lain adalah
konvergensi $(f_n)$ ke $g$ terhadap metrik seragam.
\end{defn}

Ada banyak cara barisan fungsi dapat konvergen. Boleh dikatakan bahwa dua ragam
konvergensi yang paling umum adalah konvergensi seragam, yang baru saja kita bahas, dan
konvergensi titik demi titik.

\begin{defn} Misalkan $S$ suatu himpunan, $V$ suatu ruang linear bernorma, dan $(f_n)$
suatu barisan dalam $\fml F(S,V)$. Jika terdapat fungsi $g$ sedemikian sehingga
  \[ f_n(x) \sto g(x) \qquad \text{untuk semua $x \in S$}, \]
maka $(f_n)$
 \index{titik demi titik!konvergensi}%
 \index{konvergensi!titik demi titik}%
\df{konvergen titik demi titik} ke~$g$. Dalam hal ini kita menulis
  \[ f_n \sto g \text{\,(ptws)}. \]
Fungsi $g$ adalah
 \index{titik demi titik!limit}%
 \index{limit!titik demi titik}%
 \index{<arrowto@$f_n \sto g$ (ptws) (konvergensi titik demi titik)}%
\df{limit titik demi titik} dari fungsi-fungsi~$f_n$.

Jika $(f_n)$ merupakan barisan naik dari fungsi-fungsi bernilai real (atau real
diperluas) dan $f_n \sto g \text{\,(ptws)}$, kita menulis
 \index{<arrowmonoup@$f_n \uparrow g$ (ptws) (konvergensi monoton titik demi titik)}%
$f_n \uparrow g \text{\,(ptws)}$. Dan jika $(f_n)$ menurun serta memiliki $g$ sebagai
limit titik demi titik, kita menulis
 \index{<arrowmonodown@$f_n \downarrow g$ (ptws) (konvergensi monoton titik demi titik)}%
$f_n \downarrow g \text{\,(ptws)}$.
\end{defn}

Hubungan paling jelas antara kedua jenis konvergensi ini, yang dikenal dari setiap mata
kuliah analisis real, adalah bahwa konvergensi seragam mengakibatkan konvergensi titik demi
titik.

\begin{prop}\label{C023141} Misalkan $S$ suatu himpunan dan $V$ suatu ruang linear bernorma.
Jika barisan $(f_n)$ dalam $\fml F(S,V)$ konvergen seragam ke suatu fungsi $g$ dalam
$\fml F(S,V)$, maka $(f_n)$ konvergen titik demi titik ke~$g$.
\end{prop}

Kebalikannya tidak benar.

\begin{exam} Untuk setiap $n \in \N$, misalkan
$f_n\colon [0,1] \sto \R\colon x \mapsto x^n$. Maka barisan $(f_n)$ konvergen titik demi
titik pada $[0,1]$, tetapi tidak konvergen seragam.
\end{exam}

\begin{exam} Untuk setiap $n \in \N$ dan setiap $x \in \R^+$, misalkan
$f_n(x) = \frac1n\, x$. Maka pada setiap interval berbentuk $[0,a]$ dengan $a > 0$,
barisan $(f_n)$ konvergen seragam ke fungsi konstan~$0$. Pada interval $[0,\infty)$,
barisan itu konvergen titik demi titik ke $0$, tetapi tidak konvergen seragam.
\end{exam}

Ingat pula dari analisis real bahwa limit seragam dari fungsi-fungsi terbatas juga harus
terbatas.

\begin{prop}\label{C023151} Misalkan $S$ suatu himpunan, $V$ suatu ruang linear bernorma,
$(f_n)$ suatu barisan dalam $\fml B(S,V)$, dan $g$ suatu anggota $\fml F(S,V)$. Jika
$f_n \sto g\text{\,(unif)}$, maka $g$ terbatas.
\end{prop}

Dan limit seragam dari fungsi-fungsi kontinu bersifat kontinu.

\begin{prop}\label{C023154} Misalkan $X$ suatu ruang topologis, $V$ suatu ruang linear
bernorma, $(f_n)$ suatu barisan fungsi kontinu bernilai $V$ pada~$X$, dan $g$ suatu anggota
$\fml F(X,V)$. Jika $f_n \sto g\text{\,(unif)}$, maka $g$ kontinu.
\end{prop}

\begin{prop}\label{C023163} Jika $x$ dan $y$ merupakan unsur dari suatu ruang vektor $V$
yang dilengkapi dengan norma $\norm{\hphantom{x}}$, maka
  \[ \bigabs{\strut\,\norm{x} - \norm{y}\,} \le \norm{x - y}\,. \]
\end{prop}

\begin{cor} Norma pada suatu ruang linear bernorma merupakan fungsi kontinu seragam.
 \index{kekontinuan!norma}%
 \index{norma!kekontinuan}%
\end{cor}

\begin{notn} Misalkan $M$ suatu ruang metrik, $a \in M$, dan $r > 0$. Kita menyatakan
$\{x \in M\colon d(x,a) < r\}$, yaitu
 \index{B@$B_r(a)$, $B_r$ (bola terbuka berjari-jari $r$)}%
 \index{terbuka!bola}%
 \index{bola!terbuka}%
bola terbuka berjari-jari $r$ yang berpusat di $a$, dengan $B_r(a)$. Demikian pula, kita
menyatakan $\{x \in M\colon d(x,a) \le r\}$, yaitu
 \index{C@$C_r(a)$, $C_r$ (bola tertutup berjari-jari $r$)}%
 \index{tertutup!bola}%
 \index{bola!tertutup}%
\df{bola tertutup} berjari-jari $r$ yang berpusat di $a$, dengan $C_r(a)$, dan
$\{x \in M\colon d(x,a) = r\}$, yaitu
 \index{S@$S_r(a)$, $S_r$ (sfer berjari-jari $r$)}%
 \index{sfer}%
\df{sfer} berjari-jari $r$ yang berpusat di $a$, dengan $S_r(a)$. Jika ruang metrik tersebut
memiliki titik yang ditandai sebagai nol (khususnya, jika ruang itu merupakan ruang linear bernorma), kita menggunakan
$B_r$, $C_r$, dan $S_r$ masing-masing sebagai singkatan untuk $B_r(0)$, $C_r(0)$, dan $S_r(0)$.
\end{notn}

\begin{prop} Jika $V$ merupakan ruang linear bernorma, $x \in V$, dan $r$, $s > 0$, maka
 \begin{enumerate}[label=\rm{(\alph*)}]
   \item $B_r(\vc 0) = -B_r(\vc 0)$;
   \item $B_{rs}(\vc 0) = rB_s(\vc 0)$;
   \item $x + B_r(\vc 0) = B_r(x)$; dan
   \item $B_r(\vc 0) + B_r(\vc 0) = 2\,B_r(\vc 0)$. (Apakah secara umum
$A + A = 2A$ apabila $A$ merupakan subhimpunan dari suatu ruang vektor?)
 \end{enumerate}
\end{prop}

\begin{defn}\label{C017314} Jika $a$ dan $b$ merupakan vektor dalam ruang vektor $V$, maka
 \index{tertutup!ruas}%
 \index{ruas, tertutup}%
\df{ruas tertutup} antara $a$ dan $b$, yang dinotasikan
 \index{<bracsqr@$[a,b]$ (ruas tertutup dalam ruang vektor)}%
dengan~$[a,b]$, adalah $\{(1 - t)a + tb \colon 0 \le t \le 1\}$.
\end{defn}

\begin{cau}\label{C017317} Perhatikan bahwa terdapat sedikit pertentangan antara notasi
untuk \emph{ruas} tertutup ini, ketika diterapkan pada ruang vektor bilangan real $\R$,
dan notasi biasa untuk \emph{interval} tertutup dalam~$\R$. Dalam $\R$, ruas tertutup
$[a,b]$ sama dengan interval tertutup $[a,b]$ asalkan $a \le b$. Namun, jika $a > b$,
ruas tertutup $[a,b]$ sama dengan ruas $[b,a]$ dan memuat semua bilangan $c$ sedemikian
sehingga $b \le c \le a$, sedangkan interval tertutup $[a,b]$ kosong.
\end{cau}

\begin{defn} Subhimpunan $C$ dari suatu ruang vektor $V$ disebut
 \index{konveks!himpunan}%
\df{konveks} jika ruas tertutup $[a,b]$ termuat dalam $C$ setiap kali $a$, $b \in C$.
\end{defn}

\begin{exam}\label{C023169} Dalam ruang linear bernorma, setiap bola terbuka merupakan
himpunan konveks. Demikian pula setiap bola tertutup.
\end{exam}

\begin{prop}\label{prop_intrsctn_convex} Irisan suatu keluarga subhimpunan konveks dari
ruang vektor bersifat konveks.
\end{prop}

\begin{defn} Misalkan $V$ suatu ruang vektor. Ingat bahwa suatu \emph{kombinasi linear}
dari himpunan hingga $\{x_1, \dots, x_n\}$ yang terdiri atas vektor-vektor dalam $V$ adalah
vektor berbentuk $\sum_{k=1}^n \alpha_k x_k$, dengan
$\alpha_1, \dots, \alpha_n \in \R$. Jika
$\alpha_1 = \alpha_2 = \dots = \alpha_n = 0$, maka kombinasi linear tersebut
\emph{trivial}; jika setidaknya satu $\alpha_k$ tidak sama dengan nol, kombinasi linear
tersebut \emph{nontrivial}. Suatu kombinasi linear
$\sum_{k=1}^n\alpha_k x_k$ dari vektor-vektor $x_1, \dots, x_n$ merupakan
 \index{konveks!kombinasi}%
 \index{kombinasi!konveks}%
\df{kombinasi konveks} jika $\alpha_k \ge 0$ untuk setiap $k$ ($1 \le k \le n$) dan jika
$\sum_{k=1}^n \alpha_k = 1$.
\end{defn}

\begin{defn}\label{smplx005def} Misalkan $A$ suatu subhimpunan tak kosong dari ruang
vektor~$V$. Kita mendefinisikan
 \index{konveks!selubung}%
 \index{selubung!konveks}%
 \index{co@$\co A$ (selubung konveks dari~$A$)}%
\df{selubung konveks} dari $A$, yang dinotasikan dengan $\co A$, sebagai subhimpunan
konveks terkecil dari $V$ yang memuat~$A$.
\end{defn}

Definisi sebelumnya seharusnya langsung menimbulkan kecurigaan. Bagaimana kita tahu bahwa
\emph{ada} subhimpunan konveks ``terkecil'' dari~$V$ yang memuat~$A$? Salah satu calon yang
cukup jelas (mari kita sebut calon ``abstrak'') untuk himpunan tersebut adalah irisan dari
\emph{semua} subhimpunan konveks dari $V$ yang memuat~$A$. (Tentu saja, ini juga mungkin tidak
masuk akal---\emph{kecuali} kita tahu bahwa terdapat setidaknya satu himpunan konveks dalam
$V$ yang memuat~$A$ dan bahwa irisan keluarga semua himpunan konveks yang memuat~$A$ itu
sendiri merupakan himpunan konveks yang memuat~$A$.) Calon masuk akal lainnya (yang akan
kita sebut calon ``konstruktif'') adalah himpunan semua kombinasi konveks dari
unsur-unsur~$A$. (Namun, \emph{apakah} ini suatu himpunan konveks? Dan apakah ini himpunan
terkecil yang memuat~$A$?) Akhirnya, sekalipun kedua pendekatan ini masuk akal, apakah
keduanya merupakan definisi yang ekuivalen untuk \emph{selubung konveks}? Dengan kata
lain, apakah keduanya mendefinisikan himpunan yang sama, dan apakah himpunan itu memenuhi
definisi~\ref{smplx005def}?

\begin{prop}\label{smplx006} Definisi-definisi yang disarankan dalam paragraf sebelumnya,
yang diberi nama ``abstrak'' dan ``konstruktif'', masuk akal dan ekuivalen. Selain itu,
himpunan yang dideskripsikan oleh keduanya memenuhi definisi \emph{selubung konveks} yang diberikan
dalam~\ref{smplx005def}.
\end{prop}

\begin{prop} Jika $T\colon V \sto W$ merupakan pemetaan linear antara ruang-ruang vektor
dan $C$ merupakan subhimpunan konveks dari $V$, maka $T^{\sto}(C)$ merupakan subhimpunan
konveks dari~$W$.
\end{prop}

\begin{prop} Dalam ruang linear bernorma, tutupan setiap himpunan konveks bersifat
konveks.
\end{prop}

\begin{prop} Misalkan $V$ suatu ruang linear bernorma. Untuk setiap $a \in V$, pemetaan
$T_a\colon V \sto V\colon x \mapsto x + a$ (yang disebut
 \index{translasi}%
\df{translasi} oleh~$a$) merupakan homeomorfisme.
\end{prop}

\begin{cor} Jika $U$ merupakan subhimpunan terbuka tak kosong dari ruang linear bernorma
$V$, maka $U - U$ memuat suatu lingkungan dari~$0$.
\end{cor}

\begin{prop} Jika $(x_n)$ merupakan barisan dalam ruang linear bernorma dan $x_n \sto a$,
maka $\D\frac1n\sum_{k=1}^nx_k~\sto~a$.
\end{prop}

Dalam proposisi~\ref{00015025}, kita menunjukkan bagaimana suatu hasil kali dalam pada
ruang vektor menginduksi norma pada ruang tersebut. Wajar untuk bertanya apakah semua
norma dapat dihasilkan dengan cara ini dari suatu hasil kali dalam. Jawabannya adalah
\emph{tidak}. Proposisi berikut memberikan syarat perlu dan cukup yang sangat sederhana
agar suatu norma berasal dari hasil kali dalam.

\begin{prop} Misalkan $V$ suatu ruang linear bernorma. Terdapat hasil kali dalam pada $V$
yang menginduksi norma pada $V$ jika dan hanya jika norma tersebut memenuhi
\emph{hukum jajaran genjang}~\ref{parallelogram_law}.
\end{prop}

\begin{proof}[\emph{Petunjuk pembuktian}] Untuk membuktikan bahwa jika suatu norma memenuhi
\emph{hukum jajaran genjang}, maka norma itu diinduksi oleh hasil kali dalam, gunakan
persamaan yang diberikan dalam \emph{identitas polarisasi}
(proposisi~\ref{0001507}) sebagai definisi. Buktikan terlebih dahulu bahwa
$\langle y,x \rangle = \conj{\langle x,y \rangle}$ untuk semua $x$, $y\in V$ dan bahwa
$\langle x,x \rangle > 0$ apabila $x \ne 0$. Selanjutnya, untuk $z \in V$ sebarang,
definisikan fungsi $f\colon V \sto \C\colon x \mapsto \langle x,z \rangle$. Buktikan bahwa
   \[ f(x + y) + f(x - y) = 2f(x) \tag{$\ast$} \]
untuk semua $x$, $y \in V$. Gunakan ini untuk membuktikan bahwa
$f(\alpha x) = \alpha f(x)$ untuk semua $x \in V$ dan $\alpha \in \R$.
(Mulailah dengan $\alpha$ berupa bilangan asli.) Kemudian tunjukkan bahwa $f$ aditif. (Jika
$u$ dan $v$ merupakan unsur sebarang dari $V$, ambil $x = u + v$ dan $y = u - v$.
Gunakan~$(\ast)$.) Terakhir, buktikan bahwa $f(\alpha x) = \alpha f(x)$ untuk $\alpha$
kompleks dengan menunjukkan bahwa $f(ix) = i\,f(x)$ untuk semua $x \in V$. \ns
\end{proof}

\begin{defn} Misalkan $(X,\sfml T)$ suatu ruang topologis. Suatu keluarga
$\sfml B \subseteq \sfml T$ merupakan
 \index{basis!untuk suatu topologi}%
\df{basis} untuk $\sfml T$ jika setiap anggota $\sfml T$ merupakan gabungan anggota-anggota
$\sfml B$. Dengan kata lain, keluarga $\sfml B$ dari himpunan-himpunan terbuka merupakan
basis untuk $\sfml T$ jika untuk setiap himpunan terbuka $U$ terdapat subkeluarga
$\sfml B'$ dari $\sfml B$ sedemikian sehingga $U = \bigcup \sfml B'$.
\end{defn}

\begin{exam}\label{xmpl_top_metric} Pada bidang $\R^2$ dengan topologi biasa
(Euklides), keluarga semua bola terbuka dalam ruang tersebut merupakan basis untuk
topologinya, karena setiap himpunan terbuka dalam $\R^2$ dapat dinyatakan sebagai gabungan
bola-bola terbuka. Bahkan, dalam ruang metrik mana pun, keluarga bola terbuka merupakan
basis untuk suatu topologi pada ruang tersebut. Ini adalah
 \index{topologi!diinduksi oleh metrik}%
 \index{diinduksi!topologi}%
\df{topologi yang diinduksi oleh metrik}.
\end{exam}

Dalam praktiknya (misalnya dalam ruang metrik), sering kali lebih mudah menentukan suatu
basis untuk topologi daripada menentukan topologi itu sendiri. Namun, penting untuk
menyadari bahwa mungkin terdapat banyak basis yang berbeda untuk topologi yang sama.
Setelah suatu basis tertentu dipilih, kita menyebut anggota-anggotanya
 \index{dasar!himpunan terbuka}%
 \index{terbuka!himpunan!dasar}%
\emph{himpunan terbuka dasar}.

Kata ``menginduksi'' dipandang sebagai relasi transitif. Sebagai contoh, jika diberikan
ruang hasil kali dalam $V$, tidak seorang pun akan menyebut \emph{topologi pada $V$ yang
diinduksi oleh metrik yang diinduksi oleh norma yang diinduksi oleh hasil kali dalam}.
Orang cukup menyebut \emph{topologi pada $V$ yang diinduksi oleh hasil kali dalam}.
Kadang-kadang digunakan ``dibangkitkan oleh'' sebagai pengganti ``diinduksi oleh''.

\begin{prop} Misalkan $V$ suatu ruang vektor dengan dua norma $\norm{\,\cdot\,}_1$ dan
$\norm{\,\cdot\,}_2$. Misalkan $\sfml T_k$ topologi pada $V$ yang diinduksi oleh
$\norm{\,\cdot\,}_k$ (untuk $k = 1,2$). Jika terdapat konstanta $\alpha > 0$ sedemikian
sehingga $\norm{x}_1 \le \alpha \norm{x}_2$ untuk setiap $x \in V$, maka
$\sfml T_1 \subseteq \sfml T_2$.
\end{prop}

\begin{defn} Dua norma pada ruang vektor $V$ disebut
 \index{ekuivalen!norma}%
 \index{norma!ekuivalen}%
\df{ekuivalen} jika terdapat konstanta $\alpha$, $\beta > 0$ sedemikian sehingga untuk
semua $x \in V$
   \[ \alpha\norm{x}_1 \le \norm{x}_2 \le \beta \norm{x}_1\,. \]
\end{defn}

\begin{prop}\label{C023184} Jika $\norm{\hphantom{x}}_1$ dan
$\norm{\hphantom{x}}_2$
 \index{ekuivalen!norma!menginduksi topologi yang identik}%
merupakan norma-norma ekuivalen pada ruang vektor~$V$, maka keduanya menginduksi topologi
yang sama pada~$V$.
\end{prop}

Proposisi~\ref{C023184} memberi kita syarat cukup yang mudah diperiksa agar dua norma
menginduksi topologi yang identik pada ruang vektor. Jadi, misalnya, jika kita hendak
menunjukkan bahwa suatu subhimpunan dari ruang linear bernorma bersifat terbuka, boleh jadi
pembuktiannya dapat disederhanakan dengan mengganti norma yang diberikan dengan norma
ekuivalen.

Demikian pula, andaikan kita hendak memeriksa bahwa suatu fungsi $f$ antara dua ruang
linear bernorma bersifat kontinu. Karena kekontinuan didefinisikan dalam hal himpunan
terbuka dan norma-norma ekuivalen menghasilkan himpunan-himpunan terbuka yang persis sama
(lihat proposisi~\ref{C023184}), kita bebas mengganti norma pada domain $f$ dan kodomain
$f$ dengan norma ekuivalen mana pun yang kita kehendaki. Proses ini kadang-kadang dapat
sangat menyederhanakan argumen. (Kemungkinan penyederhanaan ini, secara kebetulan,
merupakan salah satu keuntungan utama menggunakan definisi topologis untuk kekontinuan
sejak awal.)

\begin{defn}\label{C02319} Misalkan $x = (x_n)$ suatu barisan vektor dalam ruang linear
bernorma $V$. Deret tak hingga $\sum_{k=1}^\infty x_k$ disebut
 \index{konvergen!deret}%
 \index{deret!konvergen}%
\df{konvergen} jika terdapat vektor $b \in V$ sedemikian sehingga
$\norm{b - s_n} \sto 0$ ketika $n \sto \infty$, dengan
$s_n = \sum_{k=1}^n x_k$ sebagai jumlah parsial $\text{ke-}n$ dari barisan~$x$. Deret tersebut
disebut
 \index{mutlak!konvergen!deret}%
 \index{konvergen!mutlak}%
\df{konvergen mutlak} jika $\sum_{k=1}^\infty \norm{x_k} < \infty$.
\end{defn}

\begin{exer}\label{C023191} Misalkan
 \index{l@$l_1$!sebagai ruang linear bernorma}%
 \index{l@$l_1$!barisan yang dapat dijumlahkan mutlak}%
 \index{bernorma!ruang linear!$l_1$ sebagai}%
$l_1$ merupakan himpunan semua barisan $x = (x_n)$ dari bilangan kompleks sedemikian
sehingga deret $\sum x_k$ konvergen mutlak. Jadikan $l_1$ ruang vektor dengan definisi
penjumlahan dan perkalian skalar titik demi titik yang biasa. Untuk setiap $x \in l_1$,
definisikan
 \index{<unopn@$\norm{\hphantom{x}}_1$ (norma-$1$)}%
 \index{<unopn@$\norm{\hphantom{x}}_\infty$ (norma seragam)}%
   \begin{align*}
      \norm{x}_1 &:= \sum_{k=1}^\infty \abs{x_k} \\
   \intertext{dan}
      \norm{x}_\infty &:= \sup\{\abs{x_k}\colon k \in \N\}.
   \end{align*}
(Masing-masing adalah
 \index{satu@norma-$1$!dari suatu barisan}%
 \index{norma!$1$- (dari suatu barisan)}%
norma-$1$ dan
 \index{seragam!norma!dari suatu barisan}%
\emph{norma seragam}.) Tunjukkan bahwa $\norm{\hphantom{x}}_1$ dan
$\norm{\hphantom{x}}_\infty$ keduanya merupakan norma pada~$l_1$. Kemudian buktikan atau
sangkal:
 \begin{enumerate}
  \item[(a)] Jika suatu barisan $(x^i)$ dari vektor-vektor dalam ruang linear bernorma
$\bigl(l_1, \norm{\hphantom{x}}_1\bigr)$ konvergen, maka barisan tersebut juga konvergen
dalam $\bigl(l_1, \norm{\hphantom{x}}_\infty\bigr)$.
  \item[(b)] Jika suatu barisan $(x^i)$ dari vektor-vektor dalam ruang linear bernorma
$\bigl(l_1, \norm{\hphantom{x}}_\infty\bigr)$ konvergen, maka barisan tersebut juga
konvergen dalam $\bigl(l_1, \norm{\hphantom{x}}_1\bigr)$.
 \end{enumerate}
\end{exer}








\section{Pemetaan Linear Terbatas}\label{sec_bdd_lin_maps}
Para analis bekerja dengan objek-objek yang memiliki struktur aljabar sekaligus topologis.
Pada bagian sebelumnya kita mengkaji ruang-ruang vektor yang dilengkapi dengan topologi
yang berasal dari suatu norma. Jika ruang-ruang itu merupakan objek yang sedang
dipertimbangkan, morfisme apakah yang kiranya paling menarik? Jawaban yang masuk akal,
dan memang benar, adalah \emph{pemetaan yang melestarikan struktur topologis maupun
aljabar}; yakni, pemetaan linear kontinu. Kategori yang dihasilkan dinotasikan dengan
 \index{kategori!$\cat{NLS_\infty}$ sebagai suatu}%
 \index{NLSinf@$\cat{NLS_{\boldsymbol\infty}}$!kategori}
$\cat{NLS_{\boldsymbol\infty}}$. Secara topologis, isomorfisme dalam kategori ini adalah
homeomorfisme.

Berikut adalah suatu syarat yang sangat berguna yang, untuk pemetaan linear, ternyata
ekuivalen dengan kekontinuan.

\begin{defn} Suatu pemetaan linear $T\colon V \sto W$ antara ruang-ruang linear bernorma
disebut
 \index{terbatas!pemetaan linear}%
 \index{linear!pemetaan!terbatas}%
\df{terbatas} jika $T(B)$ merupakan subhimpunan terbatas dari $W$ setiap kali $B$
merupakan subhimpunan terbatas dari~$V$. Dengan kata lain, pemetaan linear terbatas
membawa himpunan terbatas ke himpunan terbatas. Kita menyatakan
 \index{B@$\ofml B(V,W)$!pemetaan linear terbatas}%
dengan~$\ofml B(V,W)$ keluarga semua pemetaan linear terbatas dari $V$ ke~$W$. Pemetaan
linear terbatas dari suatu ruang $V$ ke dirinya sendiri sering disebut suatu
 \index{operator}%
 \index{konvensi!semua operator terbatas dan linear}%
\df{operator} (linear terbatas). Keluarga semua operator pada ruang linear bernorma $V$
dinotasikan
 \index{B@$\ofml B(V)$!operator}%
dengan~$\ofml B(V)$. Kelas ruang linear bernorma bersama pemetaan linear terbatas di antara
ruang-ruang tersebut membentuk suatu kategori. Di bawah ini kita memeriksa bahwa kategori
ini tidak lain adalah kategori~$\cat{NLS_{\boldsymbol\infty}}$.
\end{defn}

\begin{cau} Sangat penting untuk menyadari bahwa secara umum suatu pemetaan linear
terbatas \emph{bukan} fungsi terbatas dalam arti biasa dari suatu \emph{fungsi terbatas}
pada himpunan tertentu (lihat contoh~\ref{C023131} dan~\ref{C023133}). Penggunaan kata
``terbatas'' dalam kedua arti yang bertentangan ini mungkin kurang menguntungkan, tetapi
sudah mapan.
\end{cau}

\begin{exam} Pemetaan linear $T\colon \R \sto \R\colon x \mapsto 3x$ merupakan pemetaan
linear terbatas (karena memetakan subhimpunan terbatas dari $\R$ ke subhimpunan terbatas
dari~$\R$), tetapi jika dipandang hanya sebagai fungsi, $T$ tidak terbatas (karena
jangkauannya bukan subhimpunan terbatas dari~$\R$).
\end{exam}

Pengamatan berikut dapat membantu mengurangi kebingungan.

\begin{prop} Suatu pemetaan linear, kecuali jika merupakan pemetaan konstan yang membawa
setiap vektor ke nol, \emph{tidak mungkin} merupakan fungsi terbatas.
\end{prop}

\begin{defn} Misalkan $(M,d)$ dan $(N,\rho)$ ruang-ruang metrik. Suatu fungsi
$f \colon M \sto N$ disebut
 \index{seragam!kekontinuan}%
 \index{kekontinuan!seragam}%
\df{kontinu seragam} jika untuk setiap $\epsilon > 0$ terdapat $\delta > 0$ sedemikian
sehingga $\rho(f(x),f(y)) < \epsilon$ setiap kali $x$, $y \in M$ dan
$d(x,y) < \delta$.
\end{defn}

Salah satu aspek linearitas yang paling mencolok adalah bahwa untuk pemetaan linear,
konsep kekontinuan, kekontinuan di satu titik saja, dan kekontinuan seragam menyatu.
Bahkan, semua konsep tersebut persis sama dengan keterbatasan.

\begin{thm}\label{C023221} Misalkan $T \colon V \sto W$ suatu pemetaan linear antara
ruang-ruang linear bernorma. Maka pernyataan-pernyataan berikut ekuivalen:
  \begin{enumerate}[label=\rm{(\alph*)}]
    \item $T$ terbatas.
    \item Citra bola satuan tertutup di bawah $T$ merupakan himpunan terbatas.
    \item $T$ kontinu seragam pada $V$.
    \item $T$ kontinu pada $V$.
    \item $T$ kontinu di~$\vc 0$.
    \item Terdapat bilangan $M > 0$ sedemikian sehingga
$\norm{Tx} \le M\norm{x}$ untuk semua $x \in V$.
  \end{enumerate}
\end{thm}

\begin{prop} Misalkan $T\colon V \sto W$ suatu pemetaan linear terbatas antara ruang-ruang
linear bernorma dengan domain bukan ruang nol. Maka keempat bilangan berikut ada dan bernilai sama.
 \begin{enumerate}[label=\rm{(\alph*)}]
  \item $\sup\{\norm{Tx}\colon \norm{x} \le 1\}$
  \item $\sup\{\norm{Tx}\colon \norm{x} = 1\}$
  \item $\sup\{\norm{Tx}\,\norm{x}^{-1}\colon x \ne \vc 0\}$
  \item $\inf\{M > 0\colon \norm{Tx} \le M\norm{x}\ \text{ untuk semua $x \in V$}\}$
 \end{enumerate}
\end{prop}

\begin{defn}\label{defn_op_norm} Jika $T$ merupakan pemetaan linear terbatas, maka
$\norm{T}$, yang disebut
 \index{norma!pemetaan linear terbatas}%
 \index{<unopn@$\norm T$ (norma pemetaan linear terbatas~$T$)}%
\df{norma} (atau
 \index{norma!operator}%
 \index{operator!norma}%
\df{norma operator}) dari $T$, didefinisikan sebagai salah satu dari empat ungkapan dalam
proposisi sebelumnya; untuk domain nol, nilainya ditetapkan nol.
\end{defn}

\begin{exam}\label{exam_op_norm_norm} Jika $V$ dan $W$ merupakan ruang-ruang linear
bernorma, maka fungsi
   \[ \norm{\hphantom{T}}\colon \ofml B(V,W) \sto \R
      \colon T \mapsto \norm{T} \]
memang merupakan suatu norma.
\end{exam}

\begin{exam}\label{C023228} Keluarga $\ofml B(V,W)$ dari semua pemetaan linear terbatas
 \index{B@$\ofml B(V,W)$!sebagai ruang linear bernorma}%
 \index{bernorma!ruang linear!$\ofml B(V,W)$ sebagai}%
antara ruang-ruang linear bernorma itu sendiri merupakan ruang linear bernorma di bawah
norma operator yang didefinisikan di atas serta operasi penjumlahan dan perkalian skalar
titik demi titik yang biasa. Yang khususnya penting adalah keluarga
$V^* = \ofml B(V,\K)$ dari anggota-anggota \emph{kontinu} dalam~$V^{\#}$, yaitu
 \index{kontinu!fungsional linear}%
 \index{linear!fungsional!kontinu}%
 \index{fungsional!linear kontinu}%
\emph{fungsional linear kontinu}, yang juga dikenal sebagai
 \index{terbatas!fungsional linear}%
 \index{linear!fungsional!terbatas}%
 \index{fungsional!linear terbatas}%
\emph{fungsional linear terbatas}.
Ingat bahwa dalam~\ref{def_alg_dual} kita mendefinisikan dual aljabar $V^{\#}$ dari ruang
vektor~$V$. Yang lebih penting dalam pembahasan kita di sini adalah $V^*$, yang disebut
 \index{<unopurstar@$V^*$ (ruang dual dari ruang linear bernorma)}%
 \index{V@$V^*$ (ruang dual dari ruang linear bernorma)}%
 \index{dual!ruang}%
 \index{ruang!dual dari ruang linear bernorma}%
\df{ruang dual} dari~$V$. Untuk membedakannya dari dual aljabar dan dual urutan, ruang ini
kadang-kadang disebut
 \index{dual!norma}%
 \index{norma!dual}%
\df{dual norma} atau
 \index{dual!topologis}%
 \index{topologis!dual}%
\df{dual topologis} dari~$V$. Perhatikan bahwa pembedaan ini tidak diperlukan dalam ruang
berdimensi hingga karena di sana $V^* = V^{\#}$ (lihat~\ref{00015033}).
\end{exam}

\begin{prop}\label{C023228a} Keluarga $\ofml B(V,W)$ dari semua pemetaan linear terbatas
 \index{B@$\ofml B(V,W)$!sebagai ruang Banach}%
 \index{Banach!ruang!$\ofml B(V,W)$ sebagai}%
antara ruang-ruang linear bernorma bersifat lengkap (terhadap metrik yang diinduksi oleh
normanya) setiap kali $W$ lengkap. Secara khusus, ruang dual dari setiap ruang linear
bernorma bersifat lengkap.
\end{prop}
















\begin{prop}\label{0003073} Misalkan $U$, $V$, dan $W$ ruang-ruang linear bernorma. Jika
$S \in \ofml B(U,V)$ dan $T \in \ofml B(V,W)$, maka $TS \in \ofml B(U,W)$ dan
$\norm{TS} \le \norm T \norm S$.
\end{prop}

\begin{exam} Pada setiap ruang linear bernorma $V$ yang bukan ruang nol,
 \index{operator!identitas}%
 \index{identitas!operator}%
\emph{operator identitas}
   \[ \id V = I_V =I\colon V \sto V \colon v \mapsto v \]
terbatas dan $\norm{I_V} = 1$.
 \index{operator!nol}%
 \index{nol!operator}%
\emph{Operator nol}
   \[ \vc 0_V = \vc 0\colon V \sto V\colon v \mapsto \vc 0 \]
juga terbatas dan $\norm{0_V} = 0$.
\end{exam}

\begin{exam} Misalkan
$T \colon \R^2 \sto \R^3\colon (x,y) \mapsto (3x, x + 2y, x - 2y)$.
Maka~$\norm{T} = \sqrt{11}$.
\end{exam}

Banyak mahasiswa yang baru saja menyelesaikan mata kuliah kalkulus dasar merasa bahwa
diferensiasi merupakan operasi yang ``lebih baik'' daripada integrasi---mungkin
karena aturan integrasi tampak lebih rumit daripada aturan diferensiasi. Namun,
ketika kita memandang keduanya sebagai pemetaan linear, yang benar justru sebaliknya.

\begin{exam} Sebagai pemetaan linear pada ruang fungsi-fungsi yang dapat didiferensialkan
pada $[0,1]$ (dengan norma seragam), integrasi terbatas; diferensiasi tidak terbatas
(dan karena itu tidak kontinu di titik mana pun!).
\end{exam}


\begin{defn} Misalkan $f\colon V_1 \sto V_2$ suatu fungsi antara ruang-ruang linear
bernorma. Untuk $k = 1,2$, misalkan $\norm{\hphantom{M}}_k$ norma pada~$V_k$ dan $d_k$
metrik pada $V_k$ yang diinduksi oleh~$\norm{\hphantom{M}}_k$ (lihat
contoh~\ref{0001503}). Maka $f$ merupakan
 \index{isometri}%
\df{isometri} (atau
suatu pemetaan \df{isometrik}) jika $d_2(f(x),f(y)) = d_1(x,y)$ untuk semua
 $x$, $y \in V_1$. Pemetaan itu dikatakan
 \index{norma!dipertahankan}%
\df{mempertahankan norma} jika $\norm{f(x)}_2 = \norm{x}_1$ untuk semua $x \in V_1$.
\end{defn}

Sebelumnya dalam bagian ini kita telah membahas kategori
$\cat{NLS_{\boldsymbol\infty}}$ dari ruang linear bernorma dan pemetaan linear kontinu.
Dalam kategori ini, isomorfisme melestarikan struktur topologis (selain struktur ruang
vektor). Apa yang akan terjadi jika kita justru menghendaki isomorfisme melestarikan
struktur metrik dari ruang-ruang linear bernorma; yakni, menjadi isometri? Mudah
dilihat bahwa (dengan adanya linearitas) hal ini ekuivalen dengan meminta agar isomorfisme
mempertahankan norma. Dalam kategori baru ini, yang dinotasikan dengan
 \index{kategori!$\cat{NLS_{\vc 1}}$ sebagai suatu}%
 \index{NLS1@$\cat{NLS_{\vc 1}}$!kategori}
$\cat{NLS_{\vc 1}}$, morfismenya adalah
 \index{kontraktif}%
\df{pemetaan linear kontraktif}, yaitu pemetaan linear $T\colon V \sto W$ antara
ruang-ruang linear bernorma sedemikian sehingga $\norm{Tx} \le \norm x$ untuk semua
$x \in V$. Tampaknya tepat untuk menyebut $\cat{NLS_{\boldsymbol\infty}}$ sebagai
 \index{topologis!kategori}%
 \index{kategori!topologis}%
\emph{kategori topologis} ruang linear bernorma dan $\cat{NLS_{\vc 1}}$ sebagai
 \index{kategori geometris}%
 \index{kategori!geometris}%
\emph{kategori geometris} ruang linear bernorma. Banyak penulis menggunakan istilah
``isomorfisme'' tanpa pengubah untuk isomorfisme dalam kategori topologis
$\cat{NLS_{\boldsymbol\infty}}$ dan
 \index{isometrik!isomorfisme}%
 \index{isomorfisme!isometrik}%
``isomorfisme isometrik'' untuk isomorfisme dalam kategori geometris
$\cat{NLS_{\vc 1}}$. Dalam catatan ini kita akan berfokus pada kategori topologis.
Teori isometrik ruang linear bernorma, meskipun penting, agak lebih khusus.

\begin{exer} Periksalah pernyataan-pernyataan yang belum dibuktikan dalam paragraf
sebelumnya. Secara khusus, buktikan bahwa
 \begin{enumerate}
    \item[(a)] $\cat{NLS_{\boldsymbol\infty}}$ merupakan suatu kategori.
    \item[(b)] Suatu isomorfisme dalam $\cat{NLS_{\boldsymbol\infty}}$ merupakan
isomorfisme ruang vektor sekaligus homeomorfisme.
    \item[(c)] Suatu pemetaan linear antara ruang-ruang linear bernorma merupakan isometri
jika dan hanya jika pemetaan tersebut mempertahankan norma.
    \item[(d)] $\cat{NLS_{\vc 1}}$ merupakan suatu kategori.
    \item[(e)] Suatu isomorfisme dalam $\cat{NLS_{\vc 1}}$ merupakan isometri sekaligus
isomorfisme ruang vektor.
    \item[(f)] Kategori $\cat{NLS_{\vc 1}}$ merupakan subkategori dari
$\cat{NLS_{\boldsymbol\infty}}$ dalam arti bahwa setiap morfisme kategori pertama
termasuk dalam kategori kedua.
 \end{enumerate}
\end{exer}













\section{Ruang Berdimensi Hingga}
Dalam bagian ini dicantumkan beberapa fakta baku dan berguna tentang ruang berdimensi hingga. Bukti hasil-hasil ini
cukup mudah ditemukan. Lihat, misalnya, \cite{BrownP:1970}, Bagian 4.4; \cite{Conway:1990}, Bagian III.3;
\cite{Kreyszig:1978}, Bagian 2.4--5; atau \cite{Rudin:1991}, halaman 16--18.

\begin{defn} Suatu barisan $(x_n)$ dalam ruang metrik disebut
 \index{Cauchy!barisan}%
 \index{barisan!Cauchy}%
\df{barisan Cauchy} jika untuk setiap $\epsilon > 0$ terdapat $n_0 \in \N$ sedemikian sehingga
$d(x_m,x_n) < \epsilon$ apabila $m$, $n \ge n_0$. Syarat ini sering disingkat sebagai
berikut: $d(x_m, x_n) \sto 0$ ketika $m$, $n \sto \infty$ (atau $\lim_{m,n \sto
\infty}d(x_m,x_n) = 0$).
\end{defn}

\begin{defn} Suatu ruang metrik $M$ disebut
 \index{lengkap!ruang metrik}%
 \index{ruang metrik!lengkap}%
 \index{ruang!metrik lengkap}%
\df{lengkap} jika setiap barisan Cauchy dalam $M$ konvergen. (Untuk meninjau kembali fakta-fakta paling dasar tentang ruang
metrik lengkap, lihat bagian pertama Bab 20 dalam catatan daring saya~\cite{Erdman:2007}.)
\end{defn}

\begin{prop}\label{prop_fdnls_Kn} Jika $V$ adalah ruang linear bernorma berdimensi $n$ atas $\K$, maka terdapat suatu pemetaan
 \index{berdimensi hingga!ruang linear bernorma}%
dari $V$ ke $\K^n$ yang sekaligus merupakan homeomorfisme dan isomorfisme ruang vektor.
\end{prop}

\begin{cor} Setiap ruang linear bernorma berdimensi hingga bersifat lengkap.
\end{cor}

\begin{cor}\label{cor_fdvs_closed} Setiap subruang vektor berdimensi hingga dari suatu ruang linear bernorma bersifat tertutup.
\end{cor}

\begin{cor}\label{cor_fdvs_norms_equiv} Sebarang dua norma pada suatu ruang vektor berdimensi hingga bersifat ekuivalen.
\end{cor}

\begin{defn} Subhimpunan $K$ dari ruang topologis $X$ disebut
 \index{kompak}%
\df{kompak} jika setiap keluarga $\sfml U$ yang terdiri atas subhimpunan-subhimpunan terbuka dari $X$ dan gabungannya memuat $K$ memiliki
subkeluarga hingga dari $\sfml U$ yang gabungannya memuat~$K$.
\end{defn}

\begin{prop}\label{00015032} Bola satuan tertutup dalam suatu ruang linear bernorma bersifat kompak jika dan hanya jika
ruang tersebut berdimensi hingga.
\end{prop}

\begin{prop}\label{00015033} Jika $V$ dan $W$ adalah ruang-ruang linear bernorma dan $V$ berdimensi hingga, maka
setiap pemetaan linear $T \colon V \sto W$ bersifat kontinu.
\end{prop}

Misalkan $T$ suatu pemetaan linear terbatas antara ruang-ruang linear bernorma (yang mungkin berdimensi tak hingga). Sama seperti dalam kasus
berdimensi hingga, kernel dan jangkauan $T$ merupakan objek yang sangat penting. Baik dalam kasus berdimensi hingga maupun tak hingga,
keduanya merupakan subruang vektor dari ruang yang memuatnya. Selain itu, dalam kedua kasus tersebut kernel $T$ bersifat tertutup.
(Di bawah suatu fungsi kontinu, pracitra himpunan tertutup bersifat tertutup.) Akan tetapi, terdapat perbedaan besar antara
kedua kasus itu. Dalam kasus berdimensi hingga, jangkauan $T$ harus tertutup (berdasarkan~\ref{cor_fdvs_closed}). Seperti akan kita lihat
dalam Contoh~\futurexref{5.2.14}{exam_ran_nonclosed}, pemetaan linear antara ruang-ruang berdimensi tak hingga tidak harus memiliki jangkauan tertutup.

\section{Hasil Bagi Ruang Linear Bernorma}

\begin{defn}\label{quo001def} Misalkan $A$ adalah objek dalam suatu kategori konkret $\cat C$.
Suatu morfisme surjektif $A \to^\pi B$ dalam $\cat C$ disebut
 \index{hasil bagi!pemetaan!untuk kategori konkret}%
 \index{pi@$\pi$ (pemetaan hasil bagi)}%
\df{pemetaan hasil bagi} untuk $A$ jika suatu fungsi $g\colon B \sto C$ (dalam $\cat{SET}$) merupakan morfisme
(dalam $\cat C$) setiap kali $g \circ \pi$ merupakan morfisme. Suatu objek $B$ dalam $\cat C$ disebut
 \index{hasil bagi!objek}%
\df{objek hasil bagi} untuk $A$ jika merupakan citra suatu pemetaan hasil bagi untuk~$A$.
\end{defn}

\begin{exam} Dalam kategori $\cat{VEC}$ yang terdiri atas ruang-ruang vektor dan pemetaan-pemetaan linear, setiap pemetaan linear
 \index{hasil bagi!pemetaan!dalam $\cat{VEC}$}%
surjektif merupakan pemetaan hasil bagi.
\end{exam}

\begin{exam}\label{C021734} Dalam kategori $\cat{TOP}$, tidak setiap epimorfisme merupakan
 \index{hasil bagi!pemetaan!dalam $\cat{TOP}$}%
pemetaan hasil bagi.
\end{exam}

\begin{proof}[\emph{Petunjuk pembuktian}] Tinjau pemetaan identitas pada bilangan real dengan topologi yang berbeda
pada domain dan kodomainnya. \ns
\end{proof}

Butir berikut, yang semestinya sudah dikenal dari aljabar linear, menunjukkan cara menghasilkan suatu objek hasil bagi tertentu dengan
``mengambil hasil bagi oleh suatu subruang''.

\begin{defn}\label{quo002def} Misalkan $M$ adalah subruang dari ruang vektor $V$. Definisikan suatu
relasi ekuivalensi $\sim$ pada $V$ dengan
   \[ x \sim y \qquad \text{jika dan hanya jika} \qquad y - x \in M. \]
Untuk setiap $x \in V$, misalkan $[x]$ adalah kelas ekuivalensi yang memuat~$x$.
 \index{<bracsqr@$[x]$ (kelas ekuivalensi yang memuat~$x$)}%
 \index{<binop@$V/M$ (hasil bagi $V$ oleh $M$)}%
Misalkan $V/M$ adalah himpunan semua kelas ekuivalensi unsur-unsur~$V$. Untuk $[x]$ dan $[y]$
dalam $V/M$, definisikan
  \[ [x] + [y] := [x+y]\]
dan untuk $\alpha \in \K$ serta $[x] \in V/M$, definisikan
  \[\alpha[x] := [\alpha x].\]
Terhadap operasi-operasi ini, $V/M$ menjadi ruang vektor. Ruang ini adalah
 \index{hasil bagi!ruang vektor}%
 \index{ruang!hasil bagi!vektor}%
\df{ruang hasil bagi} dari $V$ oleh~$M$. Notasi $V/M$ biasanya dibaca ``$V$ modulo~$M$''.
Pemetaan linear
  \[ \pi\colon V \sto V/M\colon x \mapsto [x] \]
disebut
 \index{hasil bagi!pemetaan!untuk ruang vektor}%
 \index{pi@$\pi$ (pemetaan hasil bagi)}%
\df{pemetaan hasil bagi}.
\end{defn}

\begin{exer} Verifikasilah pernyataan-pernyataan dalam Definisi~\ref{quo002def}. Secara khusus,
tunjukkan bahwa $\sim$ merupakan relasi ekuivalensi, bahwa penjumlahan dan perkalian skalar pada
himpunan kelas-kelas ekuivalensi terdefinisi dengan baik, bahwa terhadap operasi-operasi ini $V/M$ merupakan
ruang vektor, bahwa fungsi yang di sini disebut ``pemetaan hasil bagi'' memang merupakan pemetaan hasil bagi
dalam pengertian~\ref{quo001def}, dan bahwa pemetaan hasil bagi ini bersifat linear.
\end{exer}

Hasil berikut disebut \emph{teorema dasar hasil bagi} atau
 \index{hasil bagi!teorema!untuk $\cat{VEC}$}%
 \index{teorema isomorfisme pertama}%
{teorema isomorfisme pertama} untuk ruang-ruang vektor.

\begin{thm}\label{quo005} Misalkan $V$ dan $W$ adalah ruang-ruang vektor dan $M \preceq V$. Jika
$T \in \ofml L(V,W)$ dan $\ker T \supseteq M$, maka terdapat tepat satu $\wt T \in
\ofml L(V/M\,,W)$ yang membuat diagram berikut komutatif.
 \[\xymatrix@+30pt{V \ar[d]_*+{\pi} \ar[dr]^*+{T} & \\
            V/M \ar[r]_*+{\widetilde T} & W  }\]
Selanjutnya, $\wt T$ injektif jika dan hanya jika $\ker T = M$; dan $\wt T$ surjektif
jika dan hanya jika $T$ surjektif.
\end{thm}

\begin{cor} Jika $T\colon V \sto W$ adalah pemetaan linear antara ruang-ruang vektor, maka
$\ran T \cong V/\ker T$.
\end{cor}

\begin{prop}\label{C023711} Misalkan $V$ adalah ruang linear bernorma dan $M$ adalah subruang tertutup
dari~$V$.
 \index{hasil bagi!ruang linear bernorma}%
 \index{bernorma!ruang linear!hasil bagi}%
 \index{ruang!hasil bagi!linear bernorma}%
Maka pemetaan
  \[ \norm{\hphantom{x}} \colon V/M \sto \R \colon [x]
                        \mapsto \inf\{\norm u \colon u \sim x\} \]
merupakan norma pada~$V/M$. Selanjutnya,
 \index{pi@$\pi$ (pemetaan hasil bagi)}%
 \index{hasil bagi!pemetaan!untuk ruang linear bernorma}%
pemetaan hasil bagi
  \[ \pi \colon V \sto V/M \colon x \mapsto [x] \]
merupakan surjeksi linear terbatas dengan $\norm\pi \le 1$.
\end{prop}

\begin{defn} Norma yang didefinisikan pada $V/M$ disebut
 \index{hasil bagi!norma}%
 \index{norma!hasil bagi}%
\df{norma hasil bagi} pada~$V/M$. Terhadap norma ini $V/M$ merupakan ruang linear bernorma dan disebut
 \index{ruang hasil bagi}%
 \index{ruang!hasil bagi}%
\df{ruang hasil bagi} dari $V$ oleh~$M$. Ruang ini sering disebut sebagai ``$V$ modulo $M$''.
\end{defn}

\begin{thm}[Teorema dasar hasil bagi untuk $\cat{NLS_{\boldsymbol\infty}}$]\label{C023717}
 \index{hasil bagi!teorema!untuk $\cat{NLS_{\boldsymbol\infty}}$}%
Misalkan $V$ dan $W$ adalah ruang-ruang linear bernorma dan $M$ adalah subruang tertutup dari~$V$. Jika $T$ merupakan
pemetaan linear terbatas dari $V$ ke $W$ dan $\ker T \supseteq M$, maka terdapat suatu
pemetaan linear terbatas tunggal $\wt T\colon V/M \sto W$ yang membuat diagram berikut komutatif.
  \[ \xymatrix@+30pt{V \ar[d]_*+{\pi} \ar[dr]^*+{T} & \\
             V/M \ar[r]_*+{\wt T} & W  } \]
Selanjutnya: $\norm{\wt T} = \norm{T}$; $\wt T$ injektif jika dan hanya jika $\ker T = M$;
dan $\wt T$ surjektif jika dan hanya jika $T$ surjektif.
\end{thm}

\section{Hasil Kali Ruang Linear Bernorma}\label{C0155}
Hasil kali dan koproduk, seperti hasil bagi, paling baik dideskripsikan dari segi apa yang \emph{dilakukannya}.

\begin{defn}\label{C015511} Misalkan $A_1$ dan $A_2$ adalah objek-objek dalam suatu kategori~$\cat C$. Kita mengatakan bahwa
tripel $(P,\pi_1,\pi_2)$, dengan $P$ suatu objek dan $\pi_k\colon P \sto A_k$ ($k =
1,2$) morfisme-morfisme, merupakan suatu
 \index{hasil kali!dalam suatu kategori}%
\df{hasil kali} dari $A_1$ dan $A_2$ jika untuk setiap objek $B$ dan setiap pasangan morfisme
$f_k\colon B \sto A_k$ ($k = 1,2$) terdapat morfisme tunggal $f\colon B \sto P$ sedemikian
sehingga $f_k = \pi_k \circ f$ untuk $k = 1,2$.

Sudah lazim untuk mengatakan, ``Misalkan $P$ suatu hasil kali dari \dots'' sebagai pengganti, ``Misalkan $(P,\pi_1,\pi_2)$
suatu hasil kali dari \dots''. Hasil kali $A_1$ dan $A_2$ sering ditulis sebagai $A_1 \times
A_2$ atau sebagai $\prod_{k = 1,2} A_k$.
\end{defn}

Dalam suatu kategori tertentu, hasil kali mungkin ada ataupun tidak. Merupakan fakta yang menarik dan
dasar bahwa kapan pun hasil kali itu ada, hasil kali tersebut unik (hingga isomorfisme), sehingga kita
dapat berbicara tanpa ambiguitas tentang \emph{hasil kali} dua objek. Ketika kita mengatakan bahwa suatu
objek kategori yang memenuhi satu atau beberapa syarat bersifat \emph{unik hingga isomorfisme}, tentu saja kita
bermaksud bahwa sebarang dua objek yang memenuhi syarat-syarat tersebut harus isomorfik. Kita
akan sering menggunakan frasa
 \index{esensial!unik}%
 \index{konvensi!unik secara esensial berarti unik hingga isomorfisme}%
 \index{ketunggalan!esensial}%
``unik secara esensial'' untuk ``unik hingga isomorfisme''.

\begin{prop} Dalam sebarang kategori, hasil kali (jika ada) bersifat
 \index{hasil kali!ketunggalan}%
 \index{ketunggalan!hasil kali}%
unik secara esensial.
\end{prop}

\begin{exam}\label{C015517} Dalam kategori $\cat{SET}$,
 \index{hasil kali!dalam $\cat{SET}$}%
 \index{SET@$\cat{SET}$!hasil kali dalam}%
hasil kali dua himpunan $A_1$ dan $A_2$ ada dan sesungguhnya merupakan hasil kali Kartesius biasa
$A_1 \times A_2$ beserta proyeksi koordinat biasa $\pi_k\colon A_1 \times
A_2 \sto A_k \colon (a_1,a_2) \mapsto a_k$.
\end{exam}

\begin{exam}\label{C015521} Jika $V_1$ dan $V_2$ adalah ruang-ruang vektor, kita menjadikan hasil kali Kartesius
$V_1 \times V_2$ suatu ruang vektor sebagai berikut. Definisikan penjumlahan dengan
  \[ (u,v) + (w,x) := (u + w, v + x)  \]
dan perkalian skalar dengan
  \[ \alpha (u,v) := (\alpha u, \alpha v) \,. \]
Hal ini menjadikan $V_1 \times V_2$ suatu ruang vektor, dan ruang ini beserta
proyeksi koordinat biasa $\pi_k\colon V_1 \times V_2 \sto V_k \colon (v_1,v_2)
\mapsto v_k$ ($k = 1,2$) merupakan
 \index{hasil kali!dalam $\cat{VEC}$}%
 \index{VEC@$\cat{VEC}$!hasil kali dalam}%
 \index{vektor!ruang!hasil kali}%
suatu hasil kali dalam kategori $\cat{VEC}$. Ruang ini biasanya disebut
 \index{langsung!jumlah!ruang vektor}%
 \index{jumlah!langsung}%
 \index{vektor!ruang!jumlah langsung}%
 \index{<binopdirectsum@$\oplus$ (jumlah langsung)}%
\df{jumlah langsung} dari $V_1$ dan $V_2$ dan dinotasikan dengan~$V_1 \oplus V_2$.
\end{exam}

Sering kali mungkin dan dikehendaki untuk mengambil hasil kali dari sebarang keluarga objek dalam
suatu kategori. Berikut adalah perumuman Definisi~\ref{C015511}.

\begin{defn}\label{C015524} Misalkan $\bigl(A_\lambda\bigr)_{\lambda \in \Lambda}$ adalah suatu
keluarga terindeks objek-objek dalam kategori~$\cat C$. Kita mengatakan bahwa objek $P$ beserta
suatu keluarga terindeks $\bigl(\pi_\lambda\bigr)_{\lambda \in \Lambda}$ dari morfisme-morfisme
$\pi_\lambda\colon P \sto A_\lambda$ merupakan
 \index{hasil kali!dalam suatu kategori}%
\df{hasil kali} objek-objek $A_\lambda$ jika untuk setiap objek $B$ dan setiap keluarga terindeks
$\bigl(f_\lambda\bigr)_{\lambda \in \Lambda}$ dari morfisme-morfisme $f_\lambda\colon B \sto
A_\lambda$ terdapat pemetaan tunggal $g \colon B \sto P$ sedemikian sehingga $f_\lambda =
\pi_\lambda \circ g$ untuk setiap $\lambda \in \Lambda$.
\end{defn}

Suatu kategori yang di dalamnya sebarang hasil kali ada disebut
 \index{hasil kali!lengkap}%
 \index{lengkap!hasil kali}%
\df{lengkap terhadap hasil kali}. Banyak kategori yang kita jumpai dalam catatan ini lengkap
terhadap hasil kali.

\begin{defn}\label{C015531} Misalkan $\bigl(S_\lambda\bigr)_{\lambda \in \Lambda}$ adalah suatu keluarga
terindeks himpunan.
 \index{hasil kali Kartesius}%
 \index{hasil kali!Kartesius}%
\df{Hasil kali Kartesius} dari keluarga terindeks tersebut, yang dinotasikan dengan
 \index{<<@$\prod S_\lambda$ (hasil kali Kartesius)}%
$\prod_{\lambda \in \Lambda} S_\lambda$ atau cukup $\prod S_\lambda$, adalah himpunan semua
fungsi $f \colon \Lambda \sto \bigcup S_\lambda$ sedemikian sehingga $f(\lambda) \in S_\lambda$
untuk setiap $\lambda \in \Lambda$. Pemetaan-pemetaan $\pi_\lambda \colon \prod S_\lambda \sto
S_\lambda \colon f \mapsto f(\lambda)$ merupakan
 \index{koordinat!proyeksi}%
 \index{pi@$\pi_\lambda$ (proyeksi koordinat)}%
 \index{proyeksi!ke koordinat}%
\df{proyeksi koordinat} kanonik. Dalam banyak kasus, notasi $f_\lambda$ lebih mudah digunakan
daripada~$f(\lambda)$. (Lihat, misalnya, Contoh~\ref{C015537} di bawah.)
\end{defn}

\begin{exam}\label{C015534} Suatu kasus khusus yang sangat penting dari definisi sebelumnya muncul
ketika semua himpunan $S_\lambda$ identik: misalkan $S_\lambda = A$ untuk setiap $\lambda
\in \Lambda$. Dalam hal ini, hasil kali Kartesius tersebut terdiri atas semua fungsi yang memetakan
$\Lambda$ ke~$A$. Artinya, $\prod_{\lambda \in \Lambda} S_\lambda = A^\Lambda$. Perhatikan
pula bahwa dalam kasus ini proyeksi koordinat merupakan
 \index{evaluasi!pemetaan}%
 \index{fungsi!evaluasi pada suatu titik}%
\emph{pemetaan evaluasi}. Untuk setiap $\lambda$, proyeksi koordinat $\pi_\lambda$ membawa
setiap titik $f$ dalam hasil kali tersebut (yaitu, setiap fungsi $f$ dari $\Lambda$ ke~$A$) ke
$f(\lambda)$, nilainya pada~$\lambda$. Singkatnya, setiap proyeksi koordinat merupakan
pemetaan evaluasi pada suatu titik.
\end{exam}

\begin{exam}\label{C015537} Apakah $\R^n$ itu? Ruang ini hanyalah himpunan semua $n$-tupel bilangan real.
Artinya, ruang ini adalah himpunan fungsi dari $\N_n = \{1,2,\dots,n\}$ ke~$\R$. Dengan
kata lain, $\R^n$ hanyalah singkatan untuk~$\R^{\N_n}$. Dalam $\R^n$, biasanya orang menulis $x_j$ untuk
koordinat ke-$j$ dari suatu vektor $x$, bukan~$x(j)$.
\end{exam}

\begin{exam} Dalam kategori $\cat{SET}$,
 \index{hasil kali!dalam $\cat{SET}$}%
 \index{SET@$\cat{SET}$!hasil kali dalam}%
hasil kali suatu keluarga terindeks himpunan ada dan sesungguhnya merupakan hasil kali Kartesius himpunan-himpunan
tersebut beserta proyeksi koordinat kanonik.
\end{exam}

\begin{exam}\label{C015544} Jika $\bigl(V_\lambda\bigr)_{\lambda \in \Lambda}$ adalah suatu keluarga terindeks
ruang-ruang vektor, kita menjadikan hasil kali Kartesius $\prod V_\lambda$ suatu ruang
vektor sebagai berikut. Definisikan penjumlahan dan perkalian skalar secara titik demi titik: untuk $f$, $g \in
\prod V_\lambda$ dan $\alpha \in \K$,
  \[ (f + g)(\lambda) := f(\lambda) + g(\lambda)  \]
dan
  \[ (\alpha f)(\lambda) := \alpha f(\lambda) \,. \]
Hal ini menjadikan $\prod V_\lambda$ suatu ruang vektor, yang kadang-kadang disebut
 \index{langsung!hasil kali}%
 \index{hasil kali!langsung}%
\df{hasil kali langsung} dari ruang-ruang $V_\lambda$, dan ruang ini beserta proyeksi
koordinat kanonik (yang tentu saja merupakan pemetaan linear) adalah suatu
 \index{hasil kali!dalam $\cat{VEC}$}%
 \index{VEC@$\cat{VEC}$!hasil kali dalam}%
 \index{vektor!ruang!hasil kali}%
hasil kali dalam kategori $\cat{VEC}$.
\end{exam}

\begin{exam}\label{C023411} Misalkan $V$ dan $W$ adalah ruang-ruang linear bernorma. Pada hasil kali Kartesius
$V \times W$, definisikan
  \index{<unopn@$\norm{\hphantom{x}}_2$ (norma Euklides atau norma-$2$)}%
  \index{<unopn@$\norm{\hphantom{x}}_1$ (norma-$1$)}%
  \index{<unopn@$\norm{\hphantom{x}}_u$ (norma seragam)}%
 \begin{align*}
    \norm{(x,y)}_2 &= \sqrt{\norm x^2 + \norm y^2},  \\
    \norm{(x,y)}_1 &= \norm x + \norm y,  \\
   \intertext{dan}
    \norm{(x,y)}_u &= \max\{\norm x, \norm y\}\,.
 \end{align*}
Yang pertama adalah
 \index{Euklides!norma!pada hasil kali dua ruang linear bernorma}%
 \index{norma!Euklides!pada hasil kali ruang-ruang linear bernorma}%
\df{norma Euklides} (atau
 \index{dua@norma-$2$}%
 \index{norma!$2$- (pada hasil kali ruang-ruang linear bernorma)}%
\df{norma-$2$}) pada $V \times W$, yang kedua adalah
 \index{satu@norma-$1$!pada hasil kali ruang-ruang linear bernorma}%
 \index{norma!$1$- (pada hasil kali ruang-ruang linear bernorma)}%
\df{norma-$1$}, dan yang terakhir adalah
 \index{seragam!norma!pada hasil kali ruang-ruang linear bernorma}%
 \index{norma!seragam!pada hasil kali ruang-ruang linear bernorma}%
\df{norma seragam}. Memverifikasi bahwa ketiganya merupakan norma pada $V \times W$ sangat mirip dengan
argumen yang diperlukan dalam Contoh~\ref{C023125}, \ref{C023127}, dan~\ref{C023129}. Fakta bahwa
ketiganya merupakan norma-norma ekuivalen merupakan konsekuensi dari pertidaksamaan berikut dan
Proposisi~\ref{C023184}.
   \[ \norm x + \norm y \le \sqrt 2 \sqrt{\norm x^2 + \norm y^2} \le 2 \max\{\norm x,\norm y\} \le 2(\norm x + \norm y) \]
\end{exam}

\begin{conv}\label{C023414} Dalam contoh sebelumnya, kita mendefinisikan tiga norma (ekuivalen) pada
ruang hasil kali $V \times W$. Kita akan mengambil yang kedua, $\norm{\;\;}_1$, sebagai
 \index{konvensi!pemilihan norma hasil kali}%
 \index{hasil kali!norma}%
 \index{norma!hasil kali}%
\df{norma hasil kali} pada $V \times W$. Jadi, kapan pun $V$ dan $W$ merupakan ruang-ruang linear bernorma,
kecuali dinyatakan sebaliknya, kita akan memandang hasil kali $V \times W$ sebagai ruang
linear bernorma terhadap norma ini. Biasanya kita cukup menulis $\norm{(x,y)}$, bukan
$\norm{(x,y)}_1$. Hasil kali ruang-ruang linear bernorma $V$ dan $W$ biasanya dinotasikan
dengan $V \oplus W$ dan disebut
 \index{hasil kali!ruang linear bernorma}%
 \index{bernorma!ruang linear!hasil kali}%
 \index{<binopdirectsum@$\oplus$ (jumlah langsung)}%
 \index{langsung!jumlah!ruang linear bernorma}%
 \index{bernorma!ruang linear!jumlah langsung}%
\df{jumlah langsung} dari $V$ dan~$W$. (Dalam kasus khusus $V = W = \R$, metrik apakah yang
diinduksi oleh norma hasil kali pada~$\R^2$?)
\end{conv}

\begin{exam}\label{exam_prod_nls} Tunjukkan bahwa hasil kali $V \oplus W$ dari ruang-ruang linear bernorma memang merupakan suatu hasil kali
 \index{hasil kali!dalam $\cat{NLS_{\boldsymbol\infty}}$}%
 \index{NLSinf@$\cat{NLS_{\boldsymbol\infty}}$!hasil kali dalam}%
dalam kategori $\cat{NLS_{\boldsymbol\infty}}$ yang terdiri atas ruang-ruang linear bernorma dan pemetaan-pemetaan
linear terbatas.
\end{exam}

\begin{prop} Misalkan $\bigl((x_n,y_n)\bigr)$ adalah barisan dalam jumlah langsung $V \oplus W$ dari dua ruang linear
bernorma. Buktikan bahwa $\bigl((x_n,y_n)\bigr)$ konvergen ke titik $(a,b)$ dalam $V \oplus W$ jika dan hanya
jika $x_n \sto a$ dalam $V$ dan $y_n \sto b$ dalam~$W$.
\end{prop}

\begin{prop}\label{C023427} Penjumlahan merupakan operasi kontinu pada suatu ruang linear bernorma~$V$.
 \index{kekontinuan!penjumlahan}%
 \index{penjumlahan!kekontinuan}%
Artinya, pemetaan
   \[ A \colon V \oplus V \sto V \colon (x,y) \mapsto x + y \]
bersifat kontinu.
\end{prop}

\begin{exer} Berikan bukti yang \emph{sangat} singkat (tanpa $\epsilon$, $\delta$, ataupun himpunan terbuka) bahwa jika
$(x_n)$ dan $(y_n)$ merupakan barisan-barisan dalam suatu ruang linear bernorma yang masing-masing konvergen ke $a$ dan $b$,
maka $x_n + y_n \sto a + b$.
\end{exer}

\begin{prop}\label{C023434} Perkalian skalar merupakan operasi kontinu pada suatu ruang
 \index{kekontinuan!perkalian skalar}%
 \index{skalar!perkalian!kekontinuan}%
linear bernorma~$V$, dalam arti bahwa pemetaan
   \[ S \colon \K \times V \sto V \colon  (\alpha,x) \mapsto \alpha x \]
bersifat kontinu.
\end{prop}

\begin{prop} Jika $B$ dan $C$ adalah subhimpunan-subhimpunan dari suatu ruang linear bernorma dan $\alpha$ adalah skalar, maka
 \begin{enumerate}
   \item $\clo{\alpha B} = \alpha \clo B$; dan
   \item $\clo B + \clo C  \subseteq  \clo{B + C}$.
 \end{enumerate}
\end{prop}

\begin{exam} Jika $B$ dan $C$ adalah subhimpunan-subhimpunan tertutup dari suatu ruang linear bernorma, tidak
selalu berlaku bahwa $B + C$ tertutup (dan karena itu $\clo B + \clo C$ dan $\clo{B +
C}$ tidak harus sama).
\end{exam}

\begin{prop} Misalkan $C_1$ dan $C_2$ adalah subhimpunan-subhimpunan kompak dari suatu ruang linear bernorma~$V$. Maka
   \begin{enumerate}
     \item $C_1 \times C_2$ adalah subhimpunan kompak dari $V \times V$, dan
     \item $C_1 + C_2$ adalah subhimpunan kompak dari~$V$.
   \end{enumerate}
\end{prop}

\begin{exam}\label{C023441} Misalkan $X$ suatu ruang topologis. Maka keluarga
$\fml C(X)$ dari semua fungsi kontinu bernilai skalar pada~$X$ merupakan suatu
 \index{C@$\fml C(X)$!sebagai ruang vektor}%
ruang vektor terhadap operasi penjumlahan dan perkalian skalar biasa yang didefinisikan titik demi titik.
\end{exam}

\begin{exam}\label{C023443} Misalkan $X$ suatu ruang topologis. Kita menotasikan
 \index{C@$\fml C_b(X)$!fungsi kontinu terbatas pada $X$}%
dengan $\fml C_b(X)$ keluarga semua fungsi kontinu terbatas bernilai skalar pada~$X$. Keluarga ini merupakan
 \index{C@$\fml C_b(X)$!sebagai ruang linear bernorma}%
 \index{bernorma!ruang linear!$\fml C_b(X)$ sebagai suatu}%
ruang linear bernorma terhadap norma seragam. Bahkan, ruang ini merupakan subruang dari~$\fml B(X)$
(lihat Contoh~\ref{C023131}).
\end{exam}

\begin{exam}\label{C023464} Jika $S$ adalah suatu himpunan dan $a \in S$, maka fungsional evaluasi
     \[ E_a\colon \fml B(S) \sto \K\colon f \mapsto f(a) \]
merupakan fungsional linear terbatas pada ruang $\fml B(S)$ dari semua fungsi terbatas bernilai
skalar pada~$S$. Jika $S$ kebetulan merupakan ruang topologis, kita juga dapat memandang $E_a$
sebagai anggota ruang dual dari $\fml C_b(S)$. (Dalam kedua kasus tersebut, berapakah~$\norm{E_a}$?)
\end{exam}

\begin{defn}\label{def_norm_alg} Misalkan $A$ adalah suatu aljabar yang padanya telah didefinisikan sebuah norma. Andaikan bahwa sebagai tambahan,
 \index{submultiplikatif}%
sifat \df{submultiplikatif}
  \[ \norm{xy} \le \norm x\,\norm y \]
dipenuhi untuk semua $x$, $y \in A$. Maka $A$ merupakan suatu
 \index{bernorma!aljabar}%
 \index{aljabar!bernorma}%
\df{aljabar bernorma}. Kita menetapkan satu syarat tambahan: jika aljabar $A$ beridentitas, maka
   \[ \norm{\vc 1} = 1\,. \]
Dalam kasus ini, $A$ merupakan
 \index{beridentitas!aljabar bernorma}%
 \index{bernorma!aljabar!beridentitas}%
 \index{aljabar!bernorma beridentitas}%
\df{aljabar bernorma beridentitas}.
\end{defn}

\begin{exam}\label{C023448} Dalam Contoh~\ref{C023131}, telah kita tunjukkan bahwa keluarga
 \index{B@$\fml B(S)$!sebagai aljabar bernorma}%
 \index{bernorma!aljabar!$\fml B(S)$ sebagai suatu}%
$\fml B(S)$ dari fungsi-fungsi terbatas bernilai skalar pada suatu himpunan $S$ merupakan ruang linear bernorma terhadap
norma seragam. Keluarga ini juga merupakan aljabar bernorma komutatif beridentitas.
\end{exam}

\begin{prop} Perkalian merupakan operasi kontinu pada suatu aljabar bernorma~$V$, dalam arti
bahwa pemetaan
   \[ M \colon V \times V \sto V \colon (x,y) \mapsto  xy \]
bersifat kontinu.
\end{prop}

\begin{proof}[\emph{Petunjuk pembuktian}] Jika Anda dapat membuktikan bahwa perkalian pada bilangan real
bersifat kontinu, hampir pasti Anda juga dapat membuktikan bahwa perkalian bersifat kontinu pada sebarang aljabar
bernorma. \ns
\end{proof}

\begin{exam}\label{C023454} Keluarga $\fml C_b(X)$ dari fungsi-fungsi kontinu terbatas bernilai skalar
 \index{C@$\fml C_b(X)$!sebagai aljabar bernorma}%
 \index{bernorma!aljabar!$\fml C_b(X)$ sebagai suatu}%
pada suatu ruang topologis~$X$, terhadap operasi-operasi biasa yang didefinisikan titik demi titik dan
norma seragam, merupakan aljabar bernorma komutatif beridentitas. (Lihat Contoh~\ref{C023443}.)
\end{exam}

\begin{exam}\label{C023456} Keluarga $\ofml B(V)$ dari semua operator (linear terbatas)
 \index{B@$\ofml B(V)$!sebagai aljabar bernorma beridentitas}%
 \index{bernorma!aljabar!$\ofml B(V)$ sebagai suatu}%
pada suatu ruang linear bernorma bukan nol $V$ merupakan aljabar bernorma beridentitas. (Lihat~\ref{C023228} dan~\ref{C023228a}.)
\end{exam}

\section{Koproduk Ruang Linear Bernorma}\label{C0156}
Koproduk serupa dengan hasil kali---kecuali bahwa semua anak panah dibalik.

\begin{defn}\label{C015611} Misalkan $A_1$ dan $A_2$ adalah objek-objek dalam suatu kategori~$\cat C$. Tripel
$(Q,\iota_1,\iota_2)$, dengan $Q$ suatu objek dan $\iota_k\colon A_k \sto Q$ morfisme untuk $k=1,2$, disebut
 \index{koproduk}%
\df{koproduk} dari $A_1$ dan $A_2$ jika untuk setiap objek $B$ dan setiap pasangan morfisme
$f_k\colon A_k \sto B$ (k = 1,2) terdapat morfisme tunggal $f\colon Q \sto B$ sedemikian
sehingga $f_k = f \circ \iota_k$ untuk $k = 1,2$.
\end{defn}

\begin{prop} Dalam sebarang kategori, koproduk (jika ada) bersifat
 \index{koproduk!ketunggalan}%
unik secara esensial.
\end{prop}

\begin{defn}\label{C0006127} Misalkan $A$ dan $B$ adalah himpunan-himpunan. Ketika $A$ dan $B$ saling lepas, kita
sering menggunakan notasi
 \index{<binopuniond@$\uplus$, $\biguplus$ (gabungan saling lepas)}%
$A \uplus B$ sebagai pengganti $A \cup B$ (untuk menekankan bahwa $A$ dan~$B$ saling lepas). Ketika
$C = A \uplus B$, kita mengatakan bahwa $C$ merupakan
 \index{saling lepas!gabungan}%
 \index{gabungan!saling lepas}%
\df{gabungan saling lepas} dari $A$ dan~$B$. Demikian pula, kita sering memilih untuk menulis gabungan dari
suatu keluarga himpunan $\sfml A$ yang \emph{saling lepas berpasangan} sebagai $\biguplus\sfml A$. Notasi
$\biguplus_{\lambda \in \Lambda} A_\lambda$ juga dapat digunakan untuk menyatakan gabungan dari suatu
keluarga terindeks \emph{saling lepas berpasangan} $\{A_\lambda\colon  \lambda \in \Lambda\}$ dari
himpunan-himpunan. Ketika $C = \biguplus\sfml A$ atau $C = \biguplus_{\lambda \in \Lambda} A_\lambda$, kita
mengatakan bahwa $C$ merupakan \df{gabungan saling lepas} dari keluarga yang bersesuaian.
\end{defn}

\begin{defn} Dalam definisi sebelumnya, kita memperkenalkan notasi $A \uplus B$ untuk
gabungan dua himpunan saling lepas $A$ dan~$B$. Sekarang kita memperluas notasi ini dan memperbolehkan
diri kita mengambil \emph{gabungan saling lepas} dari himpunan $A$ dan $B$ sekalipun keduanya tidak
saling lepas. Kita ``membuat keduanya saling lepas'' dengan mengidentifikasi (dengan cara yang jelas) himpunan $A$ dengan
himpunan $A' = \{(a,1)\colon a \in A\}$ dan $B$ dengan himpunan $B' = \{(b,2)\colon b \in B\}$. Kemudian kita
menotasikan gabungan $A'$ dan $B'$, yang memang saling lepas, dengan $A \uplus B$ dan menyebutnya
 \index{saling lepas!gabungan}%
 \index{gabungan!saling lepas}%
 \index{<binopuniond@$\uplus$, $\biguplus$ (gabungan saling lepas)}%
\df{gabungan saling lepas} dari $A$ dan~$B$.

Secara umum, jika $\bigl(A_\lambda\bigr)_{\lambda \in \Lambda}$ adalah suatu keluarga terindeks
himpunan,
 \index{saling lepas!gabungan}%
 \index{gabungan!saling lepas}%
\df{gabungan saling lepas}nya, \smash[b]{$\biguplus\limits_{\lambda \in \Lambda}A_\lambda$}, didefinisikan sebagai
$\{(a,\lambda)\colon \lambda \in \Lambda, a \in A_\lambda\}$.
\end{defn}

\begin{exam}\label{C015617} Dalam kategori $\cat{SET}$,
 \index{koproduk!dalam $\cat{SET}$}%
 \index{SET@$\cat{SET}$!koproduk dalam}%
koproduk dua himpunan $A_1$ dan $A_2$ ada dan merupakan gabungan saling lepas $A_1 \uplus A_2$
beserta pemetaan inklusi yang jelas $\iota_k\colon A_k  \sto A_1 \uplus A_2$ ($k = 1,2$).
\end{exam}

Menarik untuk mengamati bahwa meskipun hasil kali dan koproduk suatu koleksi hingga
objek dalam kategori $\cat{SET}$ sangat berbeda, dalam kategori $\cat{VEC}$ yang lebih kompleks
keduanya ternyata merupakan hal yang persis sama.

\begin{exam} Jika $V_1$ dan $V_2$ adalah ruang-ruang vektor, jadikan hasil kali Kartesius $V_1 \times V_2$
suatu ruang vektor seperti dalam Contoh~\ref{C015521}. Ruang ini beserta
injeksi-injeksi yang jelas merupakan suatu
 \index{koproduk!dalam $\cat{VEC}$}%
 \index{VEC@$\cat{VEC}$!koproduk dalam}%
 \index{vektor!ruang!koproduk}%
koproduk dalam kategori $\cat{VEC}$.
\end{exam}

Sering kali mungkin dan dikehendaki untuk mengambil koproduk dari sebarang keluarga objek
dalam suatu kategori. Berikut adalah perumuman Definisi~\ref{C015611}.

\begin{defn}\label{C015624} Misalkan $\bigl(A_\lambda\bigr)_{\lambda \in \Lambda}$ adalah suatu keluarga terindeks
objek-objek dalam kategori~$\cat C$. Kita mengatakan bahwa objek $C$ beserta suatu
keluarga terindeks $\bigl(\iota_\lambda\bigr)_{\lambda \in \Lambda}$ dari morfisme-morfisme
$\iota_\lambda\colon A_\lambda \sto C$ merupakan
 \index{koproduk!dalam suatu kategori}%
\df{koproduk} dari objek-objek $A_\lambda$ jika untuk setiap objek $B$ dan setiap keluarga
terindeks $\bigl(f_\lambda\bigr)_{\lambda \in \Lambda}$ dari morfisme-morfisme $f_\lambda\colon
A_\lambda \sto B$ terdapat pemetaan tunggal $g \colon C \sto B$ sedemikian sehingga $f_\lambda = g
\circ \iota_\lambda$ untuk setiap $\lambda \in \Lambda$. Notasi biasa untuk
koproduk objek-objek $A_\lambda$ adalah
 \index{<<@$\coprod A_\lambda$ (koproduk)}%
$\coprod_{\lambda \in \Lambda} A_\lambda$.
\end{defn}

\begin{exam} Dalam kategori $\cat{SET}$,
 \index{koproduk!dalam $\cat{SET}$}%
 \index{SET@$\cat{SET}$!koproduk dalam}%
koproduk suatu keluarga terindeks himpunan ada dan merupakan gabungan saling lepas himpunan-himpunan tersebut.
\end{exam}

\begin{defn} Misalkan $S$ suatu himpunan dan $V$ suatu ruang vektor.
 \index{tumpuan!fungsi}%
\df{Tumpuan} suatu fungsi $f \colon S \sto V$ adalah $\{s \in S \colon f(s) \ne \vc 0\}$.
\end{defn}

\begin{exam}\label{X_dir_sum_VEC} Jika $\bigl(V_\lambda\bigr)_{\lambda \in \Lambda}$ adalah suatu keluarga terindeks
ruang-ruang vektor, kita menjadikan hasil kali Kartesius tersebut suatu ruang vektor seperti dalam Contoh~\ref{C015544}. Himpunan fungsi $f$
dalam $\prod V_\lambda$ yang memiliki tumpuan hingga (yaitu, yang bernilai tak nol hanya pada berhingga banyak unsur domainnya) jelas
merupakan subruang dari $\prod V_\lambda$. Subruang ini adalah
 \index{langsung!jumlah}%
 \index{jumlah!langsung}%
\df{jumlah langsung} dari ruang-ruang $V_\lambda$. Ruang ini dinotasikan dengan $\bigoplus_{\lambda \in
\Lambda} V_\lambda$. Ruang ini beserta pemetaan-pemetaan injektif yang jelas (yang
bersifat linear) merupakan suatu
 \index{koproduk!dalam $\cat{VEC}$}%
 \index{VEC@$\cat{VEC}$!koproduk dalam}%
 \index{vektor!ruang!koproduk}%
koproduk dalam kategori $\cat{VEC}$.
\end{exam}

\begin{exer} Apa saja koproduk dalam kategori $\cat{NLS_{\boldsymbol\infty}}$?
\end{exer}

\begin{exer} Tentukan hasil kali dan koproduk dua ruang $V$ dan $W$ dalam kategori
$\cat{NLS_{\vc 1}}$ yang terdiri atas ruang-ruang linear bernorma dan kontraksi-kontraksi linear.
\end{exer}




















\section{Jaring}\label{C0314}
Jaring (kadang-kadang disebut \emph{barisan tergeneralisasi}) berguna untuk mencirikan banyak
sifat ruang topologis umum. Dalam ruang semacam itu, jaring pada dasarnya digunakan
dengan cara yang sama seperti barisan digunakan dalam ruang metrik.

\begin{defn} Himpunan terpraurut yang setiap pasangan elemennya memiliki batas atas disebut
 \index{terarah!himpunan}%
 \index{himpunan!terarah}%
\df{himpunan terarah}.
\end{defn}

\begin{exam}\label{C031414} Jika $S$ suatu himpunan, maka $\fin S$ merupakan himpunan terarah terhadap
inklusi $\subseteq$.
\end{exam}

\begin{defn} Misalkan $S$ suatu himpunan dan $\Lambda$ suatu himpunan terarah. Pemetaan
$x\colon \Lambda \sto S$ disebut suatu
 \index{jaring}%
\df{jaring} di $S$ (atau suatu \emph{jaring anggota-anggota $S$}). Nilai $x$ di $\lambda \in
\Lambda$ biasanya ditulis $\sbsb x\lambda$ (alih-alih $x(\lambda)$), dan jaring $x$
sendiri sering dinotasikan dengan $\bigl(\sbsb x\lambda\bigr)_{\lambda \in \Lambda}$ atau cukup
$\bigl(\sbsb x\lambda\bigr)$.
\end{defn}

\begin{exam} Contoh jaring yang paling jelas ialah barisan. Barisan merupakan fungsi yang domainnya adalah himpunan (terpraurut) $\N$ dari bilangan asli.
\end{exam}

\begin{defn} Suatu jaring $x = \bigl(\sbsb x\lambda\bigr)$ dikatakan
 \index{pada akhirnya}%
\df{pada akhirnya} memiliki suatu sifat tertentu jika terdapat $\lambda_0 \in \Lambda$ sedemikian sehingga $\sbsb x\lambda$ memiliki sifat itu
apabila $\lambda \ge \lambda_0$; dan jaring itu
 \index{sering}%
\df{sering} memiliki sifat tersebut jika untuk setiap $\lambda_0 \in \Lambda$ terdapat $\lambda \in \Lambda$ sedemikian sehingga $\lambda \ge \lambda_0$
dan $\sbsb x\lambda$ memiliki sifat itu.
\end{defn}

\begin{defn} Suatu
 \index{lingkungan}%
\df{lingkungan} dari sebuah titik dalam ruang topologis adalah sebarang himpunan yang memuat suatu himpunan terbuka yang memuat titik tersebut.
\end{defn}

\begin{defn} Suatu jaring $x$ dalam ruang topologis $X$
 \index{kekonvergenan!jaring}%
 \index{jaring!kekonvergenan}%
\df{konvergen} ke titik $a \in X$ jika jaring itu pada akhirnya berada dalam setiap lingkungan~$a$. Dalam
hal ini $a$ adalah
 \index{limit!jaring}%
 \index{jaring!limit}%
\df{limit} dari $x$, dan kita menulis
   \[ x_{{}_{\sst \lambda}} \to_{\lambda \in \Lambda}a \]
atau cukup
 \index{<arrowto@$x_\lambda \sto a$ (kekonvergenan jaring)}%
$x_{{}_{\sst \lambda}} \sto a$. Jika limit bersifat unik, kita juga dapat menggunakan notasi
$\lim_{\lambda \in \Lambda} \sbsb x\lambda = a$ atau, secara lebih sederhana, $\lim \sbsb x\lambda =
a$.

Titik $a$ di $X$ merupakan
 \index{titik gugus!jaring}%
 \index{jaring!titik gugus}%
\df{titik gugus} dari jaring $x$ jika $x$ sering berada dalam setiap lingkungan~$a$.
\end{defn}

\begin{defn} Misalkan $(X,\sfml T)$ suatu ruang topologis. Subkeluarga $\sfml S \subseteq
\sfml T$ merupakan suatu
 \index{subbasis}%
\df{subbasis} untuk topologi $\sfml T$ jika keluarga semua irisan hingga dari
anggota-anggota $\sfml S$ merupakan basis untuk~$\sfml T$.
\end{defn}

\begin{exam}\label{C031429} Misalkan $X$ suatu ruang topologis dan $a \in X$. Suatu
 \index{lingkungan!basis untuk}%
 \index{basis!lingkungan suatu titik}%
\df{basis} bagi keluarga lingkungan titik $a$ adalah keluarga $\sfml B_a$ yang terdiri atas
lingkungan-lingkungan $a$ dengan sifat bahwa setiap lingkungan $a$ memuat sekurang-kurangnya
satu anggota~$\sfml B_a$. Secara umum, terdapat banyak pilihan basis lingkungan di
suatu titik. Setelah basis semacam itu dipilih, kita menyebut anggota-anggotanya sebagai
 \index{dasar!lingkungan}%
 \index{lingkungan!dasar}%
\emph{lingkungan dasar} dari~$a$.

Misalkan $\sfml B_a$ suatu basis bagi keluarga lingkungan di $a$. Urutkan $\sfml B_a$ berdasarkan
relasi memuat (kebalikan dari inklusi); yakni, untuk $U$, $V \in \sfml B_a$, tetapkan
  \[ U \preceq V \text{ jika dan hanya jika } U \supseteq V\,. \]
Dengan demikian, $\sfml B_a$ menjadi himpunan terarah. Sekarang (dengan menggunakan aksioma pilihan) pilih satu
elemen $x_{{}_U}$ dari setiap himpunan $U \in \sfml B_a$. Maka $(\sbsb xU)$ merupakan jaring di $X$
dan $\sbsb xU \sto a$.
\end{exam}

Dua proposisi berikut menjamin bahwa agar suatu jaring konvergen ke titik $a$ dalam
ruang topologis, cukup jika jaring itu pada akhirnya berada dalam setiap lingkungan dasar (atau bahkan
lingkungan subdasar) dari~$a$.

\begin{prop}\label{C031431} Misalkan $\sfml B$ suatu basis bagi topologi pada ruang topologis
$X$ dan $a$ suatu titik di~$X$. Suatu jaring $(x_\lambda)$ konvergen ke $a$ jika dan hanya
jika jaring itu pada akhirnya berada dalam setiap lingkungan~$a$ yang termasuk dalam~$\sfml B$.
\end{prop}

\begin{prop}\label{C031434} Misalkan $\sfml S$ suatu subbasis bagi topologi pada ruang topologis
$X$ dan $a$ suatu titik di~$X$. Suatu jaring $(x_\lambda)$ konvergen ke $a$ jika dan hanya jika
jaring itu pada akhirnya berada dalam setiap lingkungan~$a$ yang termasuk dalam~$\sfml S$.
\end{prop}

\begin{exam}\label{C031437} Misalkan $J = [a,b]$ suatu interval tetap pada garis real,
$\vc x = (x_0,x_1,\dots,x_n)$ suatu $(n+1)$-tupel titik-titik di $J$, dan $\vc t =
(t_1,\dots,t_n)$ suatu $n$-tupel. Pasangan $(\vc x;\vc t)$ merupakan suatu
 \index{partisi!dengan pilihan}%
 \index{pilihan}%
\df{partisi dengan pilihan} dari interval $J$ jika:
 \begin{enumerate}
  \item $x_{k-1} < x_k$ untuk $1 \le k \le n$;
  \item $t_k \in [x_{k-1},x_k]$ untuk $1 \le k \le n$;
  \item $x_0 = a$; dan
  \item $x_n = b$.
 \end{enumerate}
Gagasannya ialah bahwa $\vc x$ \emph{mempartisi} interval tersebut menjadi subinterval-subinterval dan $t_k$ adalah
titik yang \emph{dipilih} dari subinterval $\text{ke-}k$. Jika $\vc x =
(x_0,x_1,\dots,x_n)$, maka $\{\vc x\}$ menyatakan \emph{himpunan} $\{x_0,x_1,\dots,x_n\}$.
Misalkan $P = (\vc x;\vc s)$ dan $Q = (\vc y;\vc t)$ partisi-partisi (dengan pilihan) dari
interval~$J$. Kita menulis $P \preccurlyeq Q$ dan mengatakan bahwa $Q$ merupakan suatu
 \index{penghalusan}%
\df{penghalusan} dari $P$ jika $\{\vc x\} \subseteq \{\vc y\}$. Terhadap relasi
$\preccurlyeq$, keluarga $\sfml P$ dari partisi-partisi dengan pilihan pada $J$ merupakan himpunan terarah
(tetapi \emph{bukan} terurut parsial).

Sekarang andaikan $f$ suatu fungsi terbatas pada interval $J$ dan $P = (\vc x;\vc t)$ suatu
partisi $J$ menjadi $n$ subinterval. Misalkan $\Delta x_k := x_k - x_{k-1}$ untuk $1 \le k
\le n$. Lalu definisikan
   \[ S_f(P) := \sum_{k=1}^n f(t_k)\,\Delta x_k\,. \]
Setiap jumlah $S_f(P)$ semacam itu merupakan suatu
 \index{Riemann!jumlah}%
 \index{jumlah!Riemann}%
\df{jumlah Riemann} dari $f$ pada~$J$. Perhatikan bahwa karena $\sfml P$ merupakan himpunan terarah, $S_f$ adalah suatu
jaring bilangan real. Jika jaring $S_f$ dari jumlah-jumlah Riemann konvergen, kita mengatakan bahwa fungsi
$f$
 \index{Riemann!dapat diintegralkan}%
 \index{dapat diintegralkan!secara Riemann}%
\df{dapat diintegralkan secara Riemann}. Limit jaring $S_f$ (jika ada) merupakan
 \index{Riemann!integral}%
 \index{integral!Riemann}%
\df{integral Riemann} dari $f$ pada interval~$J$ dan dinotasikan dengan $\int_a^b f(x)\,dx$.
\end{exam}

\begin{exam}\label{exam_net_indiscrete} Dalam ruang topologis $X$ dengan topologi
indiskret (yang himpunan terbukanya hanya $X$ dan himpunan kosong), setiap jaring di ruang tersebut
konvergen ke setiap titik di ruang itu.
\end{exam}

\begin{exam} Dalam ruang topologis $X$ dengan topologi diskret (yang setiap subhimpunan $X$ bersifat terbuka), suatu jaring
di ruang tersebut konvergen jika dan hanya jika jaring itu pada akhirnya konstan.
 \end{exam}

\begin{defn} Suatu ruang topologis
 \index{Hausdorff}%
 \index{pemisahan!aksioma!Hausdorff}%
 \index{topologis!ruang!Hausdorff}%
 \index{topologi!Hausdorff}%
\df{Hausdorff} adalah ruang yang setiap pasangan titik berbedanya dapat
dipisahkan oleh himpunan-himpunan terbuka; yakni, untuk setiap pasangan titik berbeda $x$ dan $y$ di ruang tersebut, terdapat
lingkungan terbuka dari $x$ dan $y$ yang saling lepas.
\end{defn}

Seperti ditunjukkan oleh Contoh~\ref{exam_net_indiscrete}, jaring dalam ruang topologis umum tidak harus memiliki
limit yang unik. Namun, suatu ruang tidak memerlukan asumsi yang terlalu membatasi untuk menyingkirkan
patologi ini. Agar limit bersifat unik, cukup disyaratkan bahwa ruang itu Hausdorff.
Menariknya, ternyata syarat ini juga perlu.

\begin{prop}\label{C031441} Jika $X$ suatu ruang topologis, maka pernyataan-pernyataan berikut ekuivalen:
 \begin{enumerate}
  \item $X$ bersifat Hausdorff.
  \item Tidak ada jaring di $X$ yang dapat konvergen ke lebih dari satu titik.
  \item Diagonal di $X \times X$ bersifat tertutup.
 \end{enumerate}
\end{prop}
(Adapun
 \index{diagonal}%
\df{diagonal} dari suatu himpunan $S$ adalah $\{(s,s) \colon s \in S\} \subseteq S \times S$.)

\begin{exam} Ingat kembali dari kalkulus dasar bahwa
 \index{parsial!jumlah}%
 \index{jumlah!parsial}%
 \index{deret tak hingga!jumlah parsial}%
 \index{deret!tak hingga}%
\df{jumlah parsial} $\text{ke-}n$ dari suatu deret tak hingga $\sum_{k=1}^\infty x_k$ didefinisikan sebagai $s_n =
\sum_{k=1}^n x_k$. Deret tersebut dikatakan
 \index{kekonvergenan!deret tak hingga}%
 \index{deret tak hingga!kekonvergenan}%
\df{konvergen} jika barisan $(s_n)$ dari jumlah-jumlah parsialnya konvergen dan, jika deret tersebut
konvergen,
 \index{jumlah!deret tak hingga}%
 \index{deret tak hingga!jumlah}%
 \index{deret!jumlah}%
\emph{jumlah}nya adalah limit barisan $(s_n)$ dari jumlah-jumlah parsial. Jadi, sebagai contoh,
deret tak hingga $\sum_{k=1}^\infty 2^{-k}$ konvergen dan jumlahnya adalah~$1$.
\end{exam}

Gagasan \emph{menjumlahkan} suatu deret tak hingga sangat bergantung pada fakta bahwa
jumlah-jumlah parsial deret tersebut terurut linear oleh inklusi. ``Jumlah-jumlah parsial'' dari suatu
jaring bilangan tidak harus terurut parsial, sehingga kita memerlukan gagasan baru tentang
``penjumlahan'', yang kadang-kadang disebut \emph{penjumlahan tak berurutan}.

\begin{notn} Jika $A$ suatu himpunan, kita notasikan
 \index{fin@$\fin A$ (keluarga subhimpunan hingga dari suatu himpunan $A$)}%
dengan $\fin A$ keluarga semua subhimpunan hingga dari~$A$. (Perhatikan bahwa keluarga ini merupakan himpunan terarah terhadap
relasi~$\subseteq$.)
\end{notn}

\begin{defn}\label{def_summable_R} Misalkan $A \subseteq \R$. Untuk setiap $F \in \fin A$, definisikan $S_F$ sebagai jumlah
bilangan-bilangan dalam~$F$, yakni, $S_F = \sum F$. Maka $S = \bigl(S_F\bigr)_{F \in \fin A}$
merupakan jaring di~$\R$. Jika jaring ini konvergen, himpunan bilangan $A$ dikatakan
 \index{dapat dijumlahkan!himpunan bilangan real}%
\df{dapat dijumlahkan}; limit jaring tersebut merupakan
 \index{jumlah!himpunan yang dapat dijumlahkan}%
\df{jumlah} dari $A$ dan dinotasikan dengan $\sum A$.
\end{defn}

\begin{defn} Suatu jaring $(x_\lambda)$ dari bilangan real disebut
 \index{menaik!jaring dalam $\R$}%
 \index{jaring!menaik}%
\df{menaik} jika $\lambda \le \mu$ mengakibatkan $x_\lambda \le x_\mu$. Suatu jaring bilangan real (atau kompleks)
disebut
 \index{terbatas!jaring dalam $\K$}%
 \index{jaring!terbatas}%
\df{terbatas} jika citranya terbatas.
\end{defn}

\begin{exam}\label{C031454} Misalkan $f(n) = 2^{-n}$ untuk setiap $n \in \N$, dan definisikan $x \colon
\fin(\N) \sto \R$ dengan $x(A) = \sum_{n \in A}f(n)$.
 \begin{enumerate}
  \item Jaring $x$ menaik.
  \item Jaring $x$ terbatas.
  \item Jaring $x$ konvergen ke 1.
  \item Himpunan bilangan $\{\frac12,\frac14,\frac18, \frac1{16},\dots\}$ dapat dijumlahkan dan
jumlahnya adalah~$1$.
  \end{enumerate}
\end{exam}

\begin{exam} Dalam ruang metrik, setiap barisan konvergen bersifat terbatas. Berikan contoh yang menunjukkan
bahwa suatu jaring konvergen dalam ruang metrik (bahkan di~$\R$) tidak harus terbatas.
\end{exam}

Dari kalkulus dasar kita telah mengenal bahwa barisan menaik yang terbatas di $\R$
bersifat konvergen. Contoh di atas merupakan kasus khusus dari pengamatan yang lebih umum: jaring
menaik yang terbatas di $\R$ bersifat konvergen.

\begin{prop}\label{C031461} Setiap jaring menaik yang terbatas
$\bigl(x_{{}_{\sst \lambda}}\bigr)_{\lambda \in \Lambda}$ dari bilangan real bersifat konvergen dan
  \[ \lim x_{{}_{\sst \lambda}} = \sup\{x_{{}_{\sst \lambda}}\colon \lambda \in \Lambda\}\,. \]
\end{prop}

\begin{exam} Misalkan $(x_k)$ suatu barisan bilangan real positif yang berbeda-beda dan $A = \{x_k \colon k \in \N\}$.
Maka $\sum A$ (sebagaimana didefinisikan di atas) sama dengan jumlah deret $\sum_{k=1}^\infty
x_k$.
\end{exam}

\begin{exam} Himpunan $\bigl\{\frac1k \colon k \in \N \bigr\}$ tidak dapat dijumlahkan dan deret tak hingga
$\sum_{k=1}^\infty \frac1k$ tidak konvergen.
\end{exam}

\begin{exam} Himpunan $\bigl\{(-1)^k \frac1k\colon k \in \N \bigr\}$ tidak dapat dijumlahkan, tetapi
deret tak hingga $\sum_{k=1}^\infty (-1)^k \frac1k$ memang konvergen.
\end{exam}

\begin{prop} Setiap subhimpunan bilangan real yang dapat dijumlahkan bersifat terhitung.
\end{prop}

\begin{proof}[\emph{Petunjuk pembuktian.}] Jika $A$ suatu himpunan bilangan real yang dapat dijumlahkan dan $n \in \N$,
berapa banyak anggota $A$ yang dapat lebih besar daripada~$1/n$? \ns
\end{proof}

Dalam Proposisi~\ref{C031441}, kita telah melihat satu contoh cara jaring
mencirikan sifat topologis: suatu ruang bersifat Hausdorff jika dan hanya jika setiap jaring
di ruang tersebut memiliki paling banyak satu limit. Berikut ini sifat topologis lain---kekontinuan di suatu titik---yang dapat dengan mudah dicirikan menggunakan
jaring.

\begin{prop}\label{C031471} Misalkan $X$ dan $Y$ ruang-ruang topologis. Suatu fungsi
 \index{kekontinuan!karakterisasi oleh jaring}%
$f\colon X \sto Y$ kontinu di titik $a$ dalam $X$ jika dan hanya jika $f\bigl(\sbsb x\lambda\bigr) \sto f(a)$ di $Y$
setiap kali $\bigl(\sbsb x\lambda\bigr)$ merupakan jaring yang konvergen ke $a$ di~$X$.
\end{prop}

Dalam ruang metrik, kita mencirikan tutupan dan interior himpunan menggunakan barisan. Untuk
ruang topologis umum, kita harus mengganti barisan dengan jaring.

Dalam mendefinisikan istilah ``interior'' dan ``tutupan'', kita memiliki pilihan serupa dengan pilihan ketika kita mendefinisikan ``selubung konveks'' dalam
\ref{smplx005def}. Terdapat definisi `abstrak' dan definisi `konstruktif'. Mungkin akan membantu jika Anda mengingat kembali
Proposisi~\ref{smplx006} dan paragraf yang mendahuluinya.

\begin{defn}\label{def_interior} Untuk suatu himpunan $A$ dalam ruang topologis $X$,
 \index{interior}%
\df{interior} himpunan tersebut adalah subhimpunan terbuka terbesar dari $X$ yang termuat dalam~$A$; yakni,
gabungan semua subhimpunan terbuka dari $X$ yang termuat dalam~$A$. Interior $A$
 \index{<unopur@$\intr A$ (interior dari~$A$)}%
dinotasikan dengan $\intr A$.
\end{defn}

\begin{prop} Definisi di atas bermakna. Lebih lanjut, suatu titik $a$ dalam ruang topologis $X$ berada dalam interior suatu subhimpunan
 \index{interior!karakterisasi oleh jaring}%
 \index{karakterisasi!oleh jaring!interior}%
 \index{jaring!karakterisasi!interior}%
$A$ dari $X$ jika dan hanya jika setiap jaring di $X$ yang konvergen ke $a$ pada akhirnya berada dalam~$A$.
\end{prop}

\begin{proof}[\emph{Petunjuk pembuktian}] Satu arah mudah. Untuk arah lainnya, gunakan Proposisi~\ref{C031431} dan jenis
jaring yang dikonstruksi dalam Contoh~\ref{C031429}. \ns
\end{proof}

\begin{defn}\label{def_closure} Untuk suatu himpunan $A$ dalam ruang topologis $X$,
 \index{tutupan}%
\df{tutupan} himpunan tersebut adalah subhimpunan tertutup terkecil dari $X$ yang memuat~$A$; yakni,
irisan semua subhimpunan tertutup dari $X$ yang memuat~$A$. Tutupan $A$
 \index{<unopu@$\clo A$ (tutupan~$A$)}%
dinotasikan dengan~$\clo A$.
\end{defn}

\begin{prop}\label{C031481} Definisi di atas bermakna. Lebih lanjut, suatu titik $b$ dalam ruang topologis $X$ termasuk dalam
 \index{tutupan!karakterisasi oleh jaring}%
 \index{karakterisasi!oleh jaring!tutupan}%
 \index{jaring!karakterisasi!tutupan}%
tutupan suatu subhimpunan $A$ dari~$X$ jika dan hanya jika terdapat suatu jaring di $A$ yang konvergen ke~$b$.
\end{prop}

\begin{cor} Suatu subhimpunan $A$ dari ruang topologis bersifat tertutup jika dan hanya jika limit setiap jaring konvergen di $A$ termasuk dalam~$A$.
\end{cor}



















\section{Teorema-Teorema Hahn--Banach}

Perlu dikatakan sejak awal bahwa merujuk pada \emph{teorema Hahn--Banach} seolah-olah hanya ada satu teorema sebenarnya agak menyesatkan.
Seorang penulis yang menyatakan sedang menggunakan \emph{teorema Hahn--Banach} mungkin menggunakan salah satu dari sekitar selusin versi
teorema tersebut atau salah satu korolarinya. Apa kesamaan berbagai versi itu? Pada dasarnya, versi-versi tersebut menjamin bahwa ruang
linear bernorma pada umumnya memiliki cukup banyak fungsional linear terbatas sehingga ruang dualnya menjadi alat yang berguna
untuk mempelajari ruang itu sendiri. Artinya, versi-versi tersebut menjamin adanya teori dualitas yang menarik dan produktif.
Di bawah ini kita membahas empat versi `teorema' tersebut dan tiga korolarinya. Kita akan mengikuti praktik yang lazim:
ketika dalam pembahasan selanjutnya kita menggunakan salah satu hasil ini, kita akan mengatakan bahwa kita menggunakan `teorema' \emph{Hahn--Banach}.

Seperti yang sering terjadi dalam matematika, pengaitan hasil-hasil ini dengan Hahn dan Banach agak ganjil. Tampaknya,
gagasan-gagasan dasarnya mula-mula berasal dari Eduard Helly. Untuk sejarah menarik tentang hasil-hasil `Hahn--Banach',
lihat~\cite{Hochstadt:1980}.

Sebagai catatan, dalam nama ``Banach'', kedua huruf ``a'' dilafalkan seperti huruf ``a'' dalam kata Indonesia \emph{bapak}, sedangkan ``ch''
kira-kira dilafalkan seperti \emph{kh}; tekanan jatuh pada suku kata pertama (BAH-nakh).

\begin{defn} Suatu fungsi bernilai real $f$ pada ruang vektor real $V$ merupakan suatu
 \index{sublinear}%
\df{fungsional sublinear} jika $f(x + y) \le f(x) + f(y)$ dan $f(\alpha x) = \alpha f(x)$
untuk semua $x$, $y \in V$ dan $\alpha \ge 0$.
\end{defn}

\begin{exam} Suatu norma pada ruang vektor real merupakan fungsional sublinear.
\end{exam}

Versi pertama \emph{teorema Hahn--Banach} kita murni merupakan aljabar linear dan hanya berlaku untuk ruang vektor \emph{real}.
Versi ini memungkinkan kita memperpanjang fungsional linear tertentu pada ruang semacam itu dari subruang ke seluruh ruang.

\begin{thm}[Teorema Hahn--Banach I]\label{HBTI} Misalkan $M$ suatu subruang dari ruang vektor real $V$ dan $p$
 \index{teorema Hahn--Banach}%
suatu fungsional sublinear pada~$V$. Jika $f$ suatu fungsional linear pada $M$ sedemikian sehingga $f \le p$ pada $M$, maka
$f$ memiliki perpanjangan $g$ ke seluruh $V$ sedemikian sehingga $g \le p$ pada~$V$.
\end{thm}

\begin{proof}[\emph{Petunjuk pembuktian}] Dalam teorema ini, notasi $h \le p$ berarti bahwa $\dom h \subseteq \dom p = V$
dan bahwa $h(x) \le p(x)$ untuk semua $x \in \dom h$. Dalam hal ini kita mengatakan bahwa $h$
 \index{didominasi}%
\df{didominasi} oleh~$p$. Kita menulis
 \index{<binrelle@$f \preccurlyeq g$ ($g$ merupakan perpanjangan~$f$)}%
$f \preccurlyeq g$ jika $g$ merupakan perpanjangan~$f$ (yakni,
jika $f(x) = g(x)$ untuk semua $x \in \dom f$).

Misalkan $\fml S$ keluarga semua fungsional linear bernilai real pada subruang-subruang $V$ yang merupakan perpanjangan~$f$ dan
didominasi oleh~$p$. Keluarga $\fml S$ terurut parsial oleh $\preccurlyeq$. Gunakan \emph{lema Zorn} untuk menghasilkan suatu
elemen maksimal $g$ dari~$\fml S$. Pembuktian selesai jika Anda dapat menunjukkan bahwa $\dom{g} = V$.

Untuk itu, andaikan demi kontradiksi bahwa terdapat vektor $c \in V$ yang tidak termasuk dalam~$D = \dom{g}$. Buktikan terlebih dahulu bahwa
   \begin{equation}\label{eq_HBTI}
      g(y) - p(y - c) \le - g(x) + p(x + c)
   \end{equation}
untuk semua $x$, $y \in D$. Simpulkan bahwa terdapat bilangan $\gamma$ yang terletak di antara supremum semua bilangan
di ruas kiri~\eqref{eq_HBTI} dan infimum semua bilangan di ruas kanan.

Misalkan $E$ rentang dari $D \cup \{c\}$. Definisikan $h$ pada $E$ dengan $h(x + \alpha c) = g(x) + \alpha\gamma$ untuk $x \in D$ dan
$\alpha \in \R$. Sekarang tunjukkan bahwa $h \in \fml S$. (Ambil $\alpha > 0$ dan tinjau secara terpisah $h(x + \alpha c)$ dan $h(x - \alpha c)$.) \ns
\end{proof}

Teorema berikutnya adalah versi Teorema~\ref{HBTI} untuk ruang dengan skalar kompleks. Teorema ini kadang-kadang disebut
 \index{teorema Bohnenblust--Sobczyk--Suhomlinov}%
\emph{teorema Bohnenblust--Sobczyk--Suhomlinov}.

\begin{thm}[Teorema Hahn--Banach II]\label{HBTII} Misalkan $V$ suatu ruang vektor kompleks, $p$ suatu seminorma pada~$V$, dan $M$ suatu subruang dari~$V$.
Jika $f\colon M \sto \C$ suatu fungsional linear sedemikian sehingga $\abs f \le p$, maka $f$ memiliki perpanjangan $g$ yang merupakan fungsional linear
pada seluruh~$V$ dan memenuhi $\abs{g} \le p$.
\end{thm}

\begin{proof}[\emph{Petunjuk pembuktian}] Hasil ini pada dasarnya merupakan korolar langsung dari versi real
\emph{teorema Hahn--Banach} yang diberikan di atas dalam~\ref{HBTI}. Misalkan $V_{\R}$ ruang vektor real yang terkait dengan~$V$ (vektornya sama,
tetapi hanya menggunakan $\R$ sebagai medan skalar). Perhatikan bahwa jika $u$ adalah bagian real dari $f$ (yakni,
 \index{R@$\Re(z)$ (bagian real dari $z \in \C$)}%
$u(x) = \Re(f(x))$ untuk setiap $x \in M$), maka $u$ merupakan fungsional $\R$-linear pada $M$ (yakni, $u(x + y) = u(x) + u(y)$ dan $u(\alpha x) = \alpha u(x)$
untuk semua $x$, $y \in M$ dan semua $\alpha \in \R$), serta $f(x) = u(x) - i\,u(ix)$. Sebaliknya, jika $u$ suatu fungsional $\R$-linear pada suatu ruang vektor,
maka $f(x) := u(x) - i\,u(ix)$ mendefinisikan suatu fungsional linear (bernilai kompleks) pada ruang tersebut. Perhatikan juga bahwa jika $u$ merupakan bagian real dari~$f$,
maka $\abs{u(x)} \le p$ untuk semua $x$ jika dan hanya jika $\abs{f(x)} \le p$ untuk semua~$x$. Dengan mengingat pengamatan-pengamatan ini, cukup terapkan~\ref{HBTI}
pada bagian real dari $f$ di~$V_\R$. \ns
\end{proof}

Berikutnya mungkin merupakan versi \emph{teorema Hahn--Banach} yang paling sering digunakan.

\begin{thm}[Teorema Hahn--Banach III]\label{HBTIII} Misalkan $M$ suatu subruang vektor dari ruang linear bernorma (real atau
 \index{teorema Hahn--Banach}%
kompleks) $V$. Maka setiap fungsional linear kontinu $f$ pada $M$ memiliki
perpanjangan $g$ ke seluruh $V$ sedemikian sehingga $\norm{g} = \norm f$.
\end{thm}

\begin{proof}[\emph{Petunjuk pembuktian}] Tinjau fungsi $p\colon V \sto \R \colon x \mapsto \norm f \norm x$. \ns
\end{proof}

Ketika membuktikan suatu teorema, kadang-kadang berguna untuk mempertimbangkan kemungkinan adanya pembuktian yang lebih ``alami'' atau bahkan lebih ``jelas''.
Sesekali kita memperoleh imbalan berupa penyederhanaan yang mencerahkan dan penjelasan yang lebih jernih atas materi yang rumit. Namun, kadang-kadang kita justru terkecoh.

\begin{exer} Kawan Anda, Fred R. Dimm, mengira bahwa ia telah menemukan pembuktian langsung yang sangat sederhana untuk
versi teorema Hahn--Banach yang diberikan dalam Teorema~\ref{HBTIII}: Suatu fungsional
linear kontinu $f$ yang didefinisikan pada subruang (linear) $M$ dari ruang linear bernorma~$V$ dapat
diperpanjang menjadi fungsional linear kontinu $g$ pada seluruh $V$ tanpa memperbesar
normanya. Menurutnya, kita hanya perlu mengambil $N$ sebagai sebarang komplemen ruang vektor
dari~$M$ dan mengambil $B$ sebagai suatu basis (Hamel) untuk $N$. Definisikan $g(e) = 0$ untuk setiap vektor $e
\in B$ dan $g = f$ pada~$M$. Kemudian perpanjang $g$ secara linear ke seluruh~$V$. Tentunya, kata Fred, $g$ linear berdasarkan konstruksinya dan normanya tidak bertambah karena kita
telah menjadikannya nol di tempat-tempat yang sebelumnya belum didefinisikan. Tunjukkan kepada Fred letak kekeliruannya.
\end{exer}

\begin{thm}[Teorema Hahn--Banach IV]\label{HBTIV} Misalkan $V$ suatu ruang linear bernorma (real atau kompleks),
 \index{teorema Hahn--Banach}%
$M$ suatu subruang vektor dari $V$, dan andaikan $y \in V$ sedemikian sehingga
 \index{dist@$\dist(y,M)$ (jarak antara $y$ dan $M$)}%
$d := \dist(y,M) > 0$. Maka terdapat fungsional $g \in V^*$ sedemikian sehingga $g^{\sto}(M) = \{0\}$,
$g(y) = d$, dan $\norm g = 1$.
\end{thm}

\begin{proof}[\emph{Petunjuk pembuktian}] Misalkan $N := \spn{(M \cup \{y\})}$ dan definisikan
$f\colon N \sto \K \colon x + \alpha y \mapsto \alpha d$ ($x \in M$, $\alpha \in \K$). \ns
\end{proof}

Kasus khusus yang paling berguna dari teorema di atas adalah ketika $M = \{\vc 0\}$ dan $y \ne \vc 0$.
Hal ini memberi tahu kita bahwa satu-satunya cara agar $f(x)$ lenyap untuk setiap $f \in V^*$ adalah jika $x$
merupakan vektor nol.

\begin{cor}\label{HBCorI} Misalkan $V$ suatu ruang linear bernorma dan $x \in V$. Jika $f(x) = 0$ untuk
setiap $f \in V^*$, maka $x = \vc 0$.
\end{cor}

\begin{defn} Suatu keluarga $\fml G(S)$ dari fungsi-fungsi bernilai skalar pada himpunan~$S$
 \index{pemisahan!pasangan titik oleh keluarga fungsi}%
\df{memisahkan} titik-titik $S$ jika untuk setiap pasangan titik berbeda $x$ dan $y$ di $S$, terdapat
fungsi $f \in \fml G(S)$ sedemikian sehingga $f(x) \ne f(y)$.
\end{defn}

\begin{cor}\label{HBCorII} Jika $V$ suatu ruang linear bernorma, maka $V^*$ memisahkan titik-titik~$V$.
\end{cor}

\begin{cor}\label{HBCorIII} Misalkan $V$ suatu ruang linear bernorma dan $\vc 0 \ne x \in V$. Maka terdapat
$f \in V^*$ sedemikian sehingga $\norm f = 1$ dan $f(x) = \norm x$.
\end{cor}

\begin{prop} Misalkan $V$ suatu ruang linear bernorma, $\{x_1,\dots,x_n\}$ suatu subhimpunan yang bebas linear
dari $V$, dan $\alpha_1, \dots, \alpha_n$ skalar-skalar. Maka terdapat $f$ dalam
$V^*$ sedemikian sehingga $f(x_k) = \alpha_k$ untuk $1\le k \le n$.
\end{prop}

\begin{proof}[\emph{Petunjuk pembuktian}] Untuk $x = \sum_{k=1}^n \beta_k x_k$, tetapkan $g(x) =
\sum_{k=1}^n \alpha_k \beta_k$. \ns
\end{proof}

\begin{prop}\label{C067131} Misalkan $V$ suatu ruang linear bernorma dan $\vc 0 \ne x \in V$. Maka
   \[ \norm x = \max\{\abs{f(x)}\colon f \in V^* \text{ dan } \norm f \le 1\}\,. \]
\end{prop}

Penggunaan $\max$ alih-alih $\sup$ dalam proposisi di atas disengaja. Proposisi tersebut menyatakan bahwa
terdapat $f$ dalam bola satuan tertutup dari $V^*$ sedemikian sehingga $f(x) = \norm x$.

\begin{defn} Suatu subhimpunan $A$ dari ruang topologis $X$ disebut
 \index{padat}%
\df{padat} di $X$ jika $\clo A = X$. Ruang $X$ disebut
 \index{separabel}%
\df{separabel} jika ruang itu memuat suatu subhimpunan padat yang terhitung.
\end{defn}

\begin{exam} Ruang $\K^n$ bersifat separabel terhadap topologi biasa; yakni,
 \index{biasa!topologi!untuk $\K^n$}%
 \index{topologi!pada $\K^n$}%
 \index{K@$\K^n$!topologi biasa pada}%
topologi yang diinduksi oleh norma biasa tersebut. (Lihat \ref{C023125}, \ref{0001503}, dan \ref{xmpl_top_metric}.)
\end{exam}

\begin{exam} Dengan topologi diskret, garis real $\R$ tidak separabel.
\end{exam}

\begin{prop} Misalkan $V$ suatu ruang linear bernorma. Jika $V^*$ separabel, maka $V$ juga separabel.
\end{prop}




















\endinput
